% CoViber — Continuous Context Replication for AI-Augmented Knowledge Work
%
% Source layout: this file assembles individual sections from paper/sections/
% and the bibliography from paper/references.bib. Build with:
%   cd paper/ && make          # produces whitepaper.pdf
% or:
%   pdflatex whitepaper && bibtex whitepaper && pdflatex whitepaper && pdflatex whitepaper
%
% First-draft status: sections are the output of the paper-drafter workflow
% (task w0wqitbq7). 4 of 12 sections passed adversarial verification cleanly;
% the other 8 have refutation notes at the top of their .tex files pointing
% at specific factual claims that need correction before submission. See
% paper/UNRESOLVED.md for the full list of open citations and eval numbers.

\documentclass[11pt,letterpaper]{article}
\usepackage[margin=1in]{geometry}
\usepackage[utf8]{inputenc}
\usepackage[T1]{fontenc}
\usepackage{lmodern}
\usepackage{microtype}

% Standard scientific-paper packages
\usepackage{amsmath, amssymb, amsthm}
\usepackage{siunitx}
\usepackage{graphicx}
\usepackage{booktabs}
\usepackage{multirow}
\usepackage{xcolor}
\usepackage{hyperref}
\usepackage[capitalize,noabbrev]{cleveref}
\usepackage{algorithm}
\usepackage{algpseudocode}
\usepackage{listings}

% Code-listing style tuned for Python
\lstset{
  basicstyle=\ttfamily\small,
  keywordstyle=\bfseries,
  commentstyle=\itshape\color{gray},
  showstringspaces=false,
  frame=single,
  frameround=tttt,
  breaklines=true,
  language=Python,
}

% Hyperref polish
\hypersetup{
  colorlinks=true,
  linkcolor=blue!50!black,
  citecolor=blue!50!black,
  urlcolor=blue!50!black,
}

\title{
  CoViber: Continuous Context Replication\\
  for AI-Augmented Knowledge Work
}
\author{Punit Mishra}
\date{\today}

\begin{document}
\maketitle

% -----------------------------------------------------------------------
%  §1  Abstract + Introduction
% -----------------------------------------------------------------------
% Verified: False
% Refutation notes:
%   Two concrete misstatements against the shipped code, both grounds for refutation under the "default/constant misstated" rule:
%   
%   (1) URGENCY SIGNAL COUNT (§1.2, "Why inference-free where possible"): The draft states the urgency function is "a weighted sum of **nine** explicit signals — `@mention`, `priority_sender`, `action_word`, `question`, `unread`, `no_reply`, `collaborator`, and three mutually exclusive age tiers". That enumeration is 7 non-age signals + 3 age tiers = **ten** signals. `coviber/urgency.py` `DEFAULT_WEIGHTS` has exactly 10 keys: mention, priority_sender, action_word, question, unread, no_reply, collaborator, age_7d, age_48h, age_24h. The paper's own §Contributions bullet on urgency later says "multi-signal urgency triage function" and the abstract says "multi-signal", but §1.2 explicitly numbers them "nine" while listing ten. This is an inconsistency the reader will catch immediately against Definition 4 later, and it directly contradicts urgency.py.
%   
%   (2) LOADER SIZE (Abstract): The abstract claims "a $\approx$10-line `Loader` that emits a canonical record". The six shipped loaders (jsonl.py 39, demo.py 71, imap.py 82, slackexport.py 95, webscrape.py 102, mbox.py 156 lines) are all an order of magnitude above ~10 lines. Even the abstract Loader base class in coviber/loaders/base.py is 41 lines. §Contributions correctly and consistently says "each in $O(10^{2})$ lines" — that lower bound in the abstract directly contradicts the same paper's later claim and the shipped code. If the intent is "the source-specific *seam* is ~10 lines" (as the base class method signature effectively is), the abstract needs to say so; taken as written it is refuted by every loader that actually ships.
%   
%   Other technical claims verified as accurate against the code:
%   - "seven MCP tools over stdio" — recall, catch_me_up, who_is, project_status, graph_summary, voice_profile, refresh (mcp_server.py: 7 @mcp.tool()-decorated functions, mcp.run() over stdio).
%   - "$U(r) \in [0,14]$" at default weights — urgency.py comment "Max achievable at defaults = 11 (non-age signals) + 3 (max age tier) = 14"; contract pinned by test_urgency_contract.
%   - "content-hashed on its immutable fields (`text`, `subject`, `from_name`, `source`)" — record.py `__post_init__` calls `_hash(self.text, self.subject, self.from_name, self.source)`.
%   - "Six built-in loaders … (demo, jsonl, mbox, imap, slackexport, webscrape)" — coviber/loaders/ ls matches exactly.
%   - "core has no runtime dependencies outside the Python standard library; `[search]`, `[scrape]`, `[qdrant]`, and `[mcp]` are opt-in extras" — pyproject.toml `dependencies = []`, optional-dependencies keys exactly {scrape, search, qdrant, mcp, all, dev}.
%   - "advisory file lock across write boundaries only, so concurrent ingest and MCP-server processes on the same data directory interleave safely" — store.py `_write_lock` uses fcntl/msvcrt on `.lock`, docstring "Readers never take it".
%   - "configuration re-read per call so that operator edits take effect without a server restart" — mcp_server.py `_settings()` docstring "Re-read per call so every tool sees current config without a server restart".
%   - "empty-corpus sentinel prevents fabricated style claims" — persona.py `VoiceProfile.system_prompt` checks `if self.n_messages == 0: return "No self-authored writing samples were found …"`.
%   - "token-boundary skip filter" — urgency.py `should_skip` uses `re.search(rf"(?<![a-z0-9]){re.escape(s)}(?![a-z0-9])", …)` and DEFAULT_FYI matches on token boundaries per audit finding L5/#16.
%   - "per-project pre-compiled word-boundary regex" — workgraph.py `self._project_patterns = [(p, re.compile(rf"\b{re.escape(p.lower())}\b")) for p in self.known_projects]`.
%   - "four entity classes (people, projects, channels, tickets)" — WorkGraph.__init__ declares exactly these four defaultdicts.
%   
%   Fix the two claims above (10 signals, not 9; drop or reframe "~10-line Loader" in the abstract) and the section is otherwise fully backed by the shipped code.
% Notes to assembler: "Section labels I referenced that the later sections must declare: \\label{sec:related} (§2), \\label{sec:system} (§3), \\label{sec:record} (§3.2), \\label{sec:ingestion} (§3.3), \\label{sec:workgraph} (§3.4), \\label{sec:urgency} (§3.5), \\label{sec:persona} (§3.6), \\label{sec:mcp} (§3.8), \\label

\begin{abstract}
Large language models forget everything between sessions. The cloud-hosted memory features shipped by leading assistants close that gap by uploading user context to a vendor, which is unacceptable for professional knowledge work under privacy, compliance, or air-gap constraints. This paper describes \emph{CoViber}, a local-first memory engine for AI-augmented knowledge work that exposes continuously-updated personal context to any Model Context Protocol (MCP) client without cloud egress. The system rests on four contributions: (i) a loader-agnostic ingestion architecture in which the only source-specific code is a $\approx$10-line \texttt{Loader} that emits a canonical record, so email, chat exports, web pages, and internal APIs flow through the same downstream pipeline; (ii) an incremental cross-source \emph{WorkGraph} that unifies people, projects, channels, and tickets in $O(1)$ per record; (iii) an inference-free statistical voice model that produces a drafting system prompt from a user's own sent messages without invoking an LLM; and (iv) a configurable multi-signal urgency function $U(r) \in [0,14]$ with a token-boundary skip filter for triage. All four pieces are exposed through seven MCP tools over stdio. The reference implementation is Apache-2.0, dependency-light (the core is stdlib-only), and evaluated on a fully synthetic corpus that ships with the package; benchmark numbers, at \SI{XXX}{records\per\second} ingest and \SI{XXX}{\milli\second} warm semantic search on $10^{5}$ records, are reproducible with \texttt{python bench/run.py}.
\end{abstract}

\section{Introduction}
\label{sec:intro}

Modern large language models (LLMs) are stateless across sessions. Each new
conversation begins with an empty context window, and any prior interaction with
the user --- what they are working on, whom they collaborate with, which
messages already received a reply, what voice they write in --- must be
re-established from scratch or reconstructed at query time from an external
store. This limitation is well-known and has motivated two broad classes of
remedies: static retrieval-augmented generation (RAG) over a corpus of
documents~\cite{TODO-rag-lewis-2020}, and vendor-hosted memory features that
persist user data in the model provider's cloud.

Neither class addresses the working conditions of a professional knowledge
worker. Static RAG assumes a slowly-changing corpus indexed once and queried
many times; the target here is the opposite --- an inbox, a chat, a code-review
queue, a scraped internal wiki --- where records arrive continuously and their
salience decays on the order of hours. Vendor-hosted memory, by contrast,
matches the continuous-update assumption but violates a hard constraint in every
regulated or IP-sensitive workplace: user context (potentially containing
customer data, unreleased architecture, or personnel matters) is transmitted to
and retained on infrastructure the user does not control. Air-gapped
environments, financial-services teams, healthcare, and any organisation with a
DLP policy simply cannot use such features~\cite{TODO-privacy-llm-2024}.

The gap between what LLMs need to be useful across sessions and what a
knowledge worker can safely feed them is the problem this paper addresses. We
argue that a memory engine for AI-augmented knowledge work should satisfy four
properties simultaneously:

\begin{enumerate}
    \item \textbf{Local-first execution.} Records, embeddings, and derived
    structures live on the user's machine; no data leaves the box unless the
    operator explicitly configures a remote store.
    \item \textbf{Loader-agnostic core.} Adding a new source (a bespoke CRM, a
    ticketing system, an internal API) should require writing a source-specific
    adapter only, not touching the graph, triage, or search components.
    \item \textbf{Inference-free enrichment where possible.} Enrichment steps
    that can be expressed statistically --- entity linking, voice modelling,
    urgency scoring --- should not require an LLM call. Fewer inference calls
    means lower latency, no per-query cost, and no additional trust boundary
    around a hosted model.
    \item \textbf{Standardised exposure to arbitrary LLM clients.} The memory
    should be consumable from any assistant a user might reach for, without a
    bespoke integration per client.
\end{enumerate}

\subsection{The design commitment}

CoViber realises these four properties with a deliberately narrow architectural
seam. All source-specific code lives behind a single \texttt{Loader} interface
whose only obligation is to yield instances of a canonical \texttt{Record}
dataclass (\S\ref{sec:record}). The record schema is content-hashed on its
immutable fields (\texttt{text}, \texttt{subject}, \texttt{from\_name},
\texttt{source}) so that re-ingesting the same message is a natural
no-op rather than a duplication event. Everything downstream --- deduplication,
the cross-source WorkGraph, the urgency queue, the semantic search index, the
voice profile, and the MCP tool surface --- is loader-agnostic. Adding a source
means writing roughly ten lines of Python; nothing in the ``brain'' changes.

The exposure layer is the Model Context Protocol (MCP)~\cite{TODO-mcp-spec}, an
emerging open standard for tools and context providers to expose capabilities
to LLM clients over a well-defined JSON-RPC schema. MCP is the right layer for a
memory engine because it decouples the memory implementation from the client:
the same seven CoViber tools --- \texttt{recall}, \texttt{catch\_me\_up},
\texttt{who\_is}, \texttt{project\_status}, \texttt{graph\_summary},
\texttt{voice\_profile}, and \texttt{refresh} --- are consumed identically by
Claude Desktop, Claude Code, and any other MCP-conformant client. The transport
in the default configuration is standard input/output to a locally-spawned
process, which keeps the entire round-trip on the machine.

\subsection{Why inference-free where possible}

Two of the four subsystems that could plausibly delegate to an LLM ---
persona/voice modelling and urgency scoring --- instead use classical statistics
and configurable rule-based signals. This is a deliberate tradeoff. LLM-based
voice modelling would produce marginally more nuanced style transfer, but at
the cost of a network round-trip (or a resident local model) per drafting
call, plus opacity: the user cannot audit why a particular draft came out the
way it did. CoViber's \texttt{VoiceProfile} is a small dataclass containing
counts and rates (average words per message, formality index, opener/closer
frequencies, signature vocabulary) computed from the user's own sent messages
(\S\ref{sec:persona}), from which a compact drafting system prompt is emitted
verbatim. Every field is inspectable; nothing is hallucinated.

The urgency function similarly refuses inference. It is a weighted sum of
nine explicit signals --- \texttt{@mention}, \texttt{priority\_sender},
\texttt{action\_word}, \texttt{question}, \texttt{unread}, \texttt{no\_reply},
\texttt{collaborator}, and three mutually exclusive age tiers --- with a
token-boundary skip filter that suppresses bots, digests, and out-of-office
noise before scoring. The default weights bound $U(r) \in [0, 14]$; the weights
themselves, the priority-sender list, the collaborator list, and the skip
lexicons are all configuration. The result is a triage that is fast (linear
scan, no inference), reproducible (same input, same score), and explainable
(the scored output carries the list of signals that fired).

\subsection{Contributions}
The specific contributions of this paper are as follows:
\begin{itemize}
    \item \textbf{A loader-agnostic ingestion architecture}
    (\S\ref{sec:ingestion}) in which sources plug in through a single
    \texttt{Loader} seam, with a content-hashed \texttt{Record} whose
    identity is invariant under re-ingest. Six built-in loaders are shipped
    (\texttt{demo}, \texttt{jsonl}, \texttt{mbox}, \texttt{imap},
    \texttt{slackexport}, \texttt{webscrape}), each in $O(10^{2})$ lines.
    \item \textbf{An incremental cross-source WorkGraph}
    (\S\ref{sec:workgraph}) that maintains four entity classes (people,
    projects, channels, tickets) with bidirectional edges built at $O(1)$
    per record. Entity keys are case-normalised; ticket identifiers are
    canonicalised; project matching is bounded by a per-project pre-compiled
    word-boundary regex.
    \item \textbf{An inference-free statistical voice model}
    (\S\ref{sec:persona}) that emits a drafting system prompt directly from
    a corpus of self-authored messages, with no LLM in the loop. An
    empty-corpus sentinel prevents fabricated style claims.
    \item \textbf{A multi-signal urgency triage function} (\S\ref{sec:urgency})
    $U(r) : \mathrm{Record} \to \mathbb{N}$ with configurable weights, a
    token-boundary skip filter to eliminate false-positive suppression on
    substring matches, and a documented contract $U(r) \in [0, 14]$ at
    default weights.
    \item \textbf{An MCP-native exposure surface} (\S\ref{sec:mcp}) of seven
    tools that make the memory queryable by any MCP-conformant LLM client
    over stdio, with configuration re-read per call so that operator edits
    take effect without a server restart.
\end{itemize}

Beyond the algorithmic contributions, the reference implementation is a design
artefact in its own right: the core package has no runtime dependencies outside
the Python standard library; \texttt{[search]}, \texttt{[scrape]},
\texttt{[qdrant]}, and \texttt{[mcp]} are opt-in extras, each degrading
gracefully when absent. Persistence uses an advisory file lock across write
boundaries only, so concurrent ingest and MCP-server processes on the same
data directory interleave safely without a transactional wrapper around the
whole pipeline. All numbers reported in \S\ref{sec:eval} are produced by a
single benchmark harness (\texttt{bench/run.py}) over a fully synthetic
\texttt{demo} corpus that ships with the package.

\subsection{Roadmap}
The rest of the paper is organised as follows. Section~\ref{sec:related}
positions CoViber against static-corpus RAG, agent frameworks, and
vendor-hosted memory. Section~\ref{sec:system} formalises the four subsystems
--- record schema and ingestion (\S\ref{sec:ingestion}), the WorkGraph
(\S\ref{sec:workgraph}), the persona engine (\S\ref{sec:persona}), the urgency
function (\S\ref{sec:urgency}), and the MCP tool surface (\S\ref{sec:mcp}) ---
with algorithms, definitions, and complexity bounds against the shipped
implementation. Section~\ref{sec:eval} reports throughput and search-latency
measurements on the synthetic corpus at multiple scales, discusses when the
default \texttt{JSONVectorStore} gives way to an optional \texttt{Qdrant}
backend, and quantifies the tradeoff. Section~\ref{sec:discussion} discusses
threats to validity, deliberate design invariants recorded in the codebase
(content-hash id boundaries, per-call config re-reads, verbatim record
storage), and directions for future work: incremental graph maintenance,
richer person-alias resolution, and community-contributed loaders.


% -----------------------------------------------------------------------
%  §2  Related Work
% -----------------------------------------------------------------------
% Verified: False
% Refutation notes:
%   The §2.4 (Voice and style modelling) claim overstates the persona feature vector. The draft says the module "computes seven scalar features --- message count, mean word count, formality score (a bounded ratio of formal to casual marker hits), emoji rate, question rate, exclamation rate", but that enumeration lists only six items, and the `VoiceProfile` dataclass in coviber/persona.py (lines 28-38) has exactly six scalar fields: `n_messages`, `avg_words`, `formality`, `emoji_rate`, `question_rate`, `exclaim_rate` (the remaining three fields — `top_openers`, `top_closers`, `vocabulary` — are lists, not scalars). Change "seven" to "six", or add a seventh scalar (e.g. include `top_openers`/`top_closers`/`vocabulary` counts in a different count and re-classify). Secondary nit worth flagging while fixing: the draft describes "the three most-common opener and closer bigrams", but the greeting/closing regexes in persona.py capture variable-length phrases (single tokens like "hi"/"hey"/"thanks" as well as bigrams like "good morning"/"thank you"), and the code takes `most_common(3)` on the raw capture-group strings, not on 2-gram statistics — "opener/closer phrases" would be accurate. Every other technical claim in the section (loader-agnostic core with a single Loader seam, 7 MCP tools listed correctly, ~180-line mcp_server.py — actually 179 lines, per-tool error strings via the try/except in refresh(), COVIBER_CONFIG re-read per call in _settings(), JSONVectorStore default, non-loopback Qdrant stderr warning via _warn_if_remote_qdrant, dedup as content hash, inference-free graph/urgency/persona, bounded linear urgency U(r) ∈ [0,14] at defaults, zero-core-deps with [mcp]/[search]/[qdrant]/[scrape] extras, empty-corpus refusal sentinel in VoiceProfile.system_prompt) is verified against the shipped code.
% Notes to assembler: Cross-section \\ref targets used: sec:arch, sec:workgraph, sec:urgency, sec:persona, sec:record, sec:mcp, sec:vectors, sec:eval, and sec:related:rag/personal/mcp/voice. The assembler must ensure the Architecture / Design section labels these anchors, or replace \\ref{...} with textual references. Pa

\section{Related Work}
\label{sec:related}

CoViber sits at the intersection of four research and engineering threads: retrieval-augmented memory for language models, personal-context replication systems, tool-augmented LLM interfaces, and computational voice modelling. In each we identify the closest neighbors and articulate what CoViber does differently. The recurring axis of difference is that CoViber treats the \emph{source of context} as the only pluggable seam (\S\ref{sec:arch}) and everything downstream --- the cross-source graph (\S\ref{sec:workgraph}), the urgency score (\S\ref{sec:urgency}), the voice profile (\S\ref{sec:persona}) --- as fixed, deterministic, and inference-free.

\subsection{Retrieval-augmented generation and memory frameworks}
\label{sec:related:rag}

The dominant approach to giving LLMs continuity is retrieval-augmented generation (RAG) over a document store, typically driven by dense-embedding similarity~\cite{TODO-rag-lewis-2020,TODO-dense-passage-retrieval}. Application-layer frameworks --- LangChain's memory abstractions~\cite{TODO-langchain-memory}, LlamaIndex's data connectors and index types~\cite{TODO-llamaindex}, and LangGraph's stateful agent graphs~\cite{TODO-langgraph} --- generalize this pattern into pipelines that ingest heterogeneous documents, chunk and embed them, and expose a retriever interface. More recently, dedicated memory services such as mem0~\cite{TODO-mem0} and Zep~\cite{TODO-zep} add long-term episodic memory with LLM-driven summarization and fact extraction on top of a vector index.

CoViber differs from this family along three axes. First, it is \emph{loader-agnostic by construction}: the only seam is the \texttt{Loader} interface (\S\ref{sec:arch}), and every downstream stage --- dedup, work-graph, urgency, semantic search, MCP tools --- operates on the canonical \texttt{Record} schema (\S\ref{sec:record}). The pipeline is a single function that takes a \texttt{Settings} dataclass and dispatches through a registry, so a new source is $\approx$10 lines of code and no core change. LangChain and LlamaIndex are also extensible via connectors, but their downstream primitives (retrievers, agents, tools) still assume a document-centric abstraction; CoViber's downstream is a \emph{work} model (people, projects, channels, tickets), not a document model. Second, CoViber's memory is \emph{not a vector index alone}. Semantic recall is one MCP tool (\texttt{recall}) among seven; the others expose a symbolic work graph (\texttt{who\_is}, \texttt{project\_status}, \texttt{graph\_summary}), a triage queue (\texttt{catch\_me\_up}), a voice profile (\texttt{voice\_profile}), and an ingest hook (\texttt{refresh}). This co-existence of vector and symbolic memory is deliberate: the two answer different questions, and neither subsumes the other. Third, CoViber is \emph{local-first and inference-free at every core stage}. Dedup is a content hash (\S\ref{sec:record}), the graph is deterministic entity extraction with regex primitives (\S\ref{sec:workgraph}), urgency is a bounded linear scoring function (\S\ref{sec:urgency}), and the persona model is closed-form statistics over token counts (\S\ref{sec:persona}). No LLM call is required to update state; the LLM only reads through MCP. This is a strict inversion of the mem0/Zep model, where an LLM is on the write path.

\subsection{Personal-context and continuous-replication systems}
\label{sec:related:personal}

The closest neighbors are systems that model a specific user's context continuously rather than answering document queries. Rewind~\cite{TODO-rewind} records the desktop stream (screen, audio) and indexes it locally for later semantic search; Mem~\cite{TODO-mem-ai} builds a personal knowledge graph from notes and messages with LLM-driven auto-organization; Karpathy's persona notes~\cite{TODO-karpathy-persona} sketch a lightweight, statistically-informed style profile embedded in a system prompt. Enterprise variants include Glean~\cite{TODO-glean} and Microsoft's Copilot memory~\cite{TODO-copilot-memory}, both of which operate as cloud-hosted work-graphs over corporate SSO'd content.

CoViber is architecturally the same shape as Rewind and Mem --- a continuously-updated personal-context store queried through a chat interface --- and its persona component is the direct spiritual descendant of Karpathy's sketch. It differs in three ways. (i) The transport is standardized: CoViber ships an MCP server (\S\ref{sec:mcp}) whose seven tools any conformant client can call, rather than being married to a proprietary desktop app. (ii) The data path is verifiably local: records, work-graph, and vectors live on disk under a single \texttt{data\_dir}, the default vector backend is a plain JSON file (\texttt{JSONVectorStore}, \S\ref{sec:vectors}), and the config parser deliberately refuses shell-style variable expansion so credentials cannot be smuggled into a shared YAML. When a non-loopback Qdrant URL is configured, the store prints a stderr warning at startup so a copy-pasted managed-cluster URL does not silently exfiltrate embeddings. (iii) No LLM is required to run: the core has zero non-stdlib dependencies, and the optional extras (\texttt{[mcp]}, \texttt{[search]}, \texttt{[qdrant]}, \texttt{[scrape]}) each degrade gracefully when absent. In contrast, both Mem and Copilot's memory require a cloud LLM in the write path; Rewind requires its proprietary indexer. The evaluation corpus in this paper is 100\% synthetic (see \S\ref{sec:eval} and the datasheet), which is not merely a legal convenience: it establishes that continuous context replication does not require personal data to be described or benchmarked.

\subsection{MCP and tool-augmented LLMs}
\label{sec:related:mcp}

Tool-augmented language models --- ReAct~\cite{TODO-react-yao-2022}, Toolformer~\cite{TODO-toolformer}, and their descendants --- established that LLMs benefit from structured external calls beyond text generation. The Model Context Protocol (MCP)~\cite{TODO-mcp-spec} generalises this into a client--server contract: any host application (Claude Desktop, Claude Code, custom clients) can discover and invoke tools exposed by any conformant server, with schema and permission negotiated at connection time. Existing MCP servers largely fall into two classes: thin adapters over a single API (Slack, Linear, GitHub) that expose CRUD-style tools, and heavier retrieval servers that wrap a hosted vector index.

CoViber occupies a third position: it is an MCP server whose tools expose \emph{derived} local state --- a work graph, a triage queue, a voice model --- rather than raw API endpoints or a black-box index. The server is a $\approx$180-line FastMCP module (\texttt{coviber/mcp\_server.py}) that re-reads its \texttt{Settings} on every tool call, so a \texttt{COVIBER\_CONFIG} edit takes effect without a restart, and any per-tool failure (unknown loader, missing file, corrupt work-graph JSON) is returned as a readable string rather than propagating a Python traceback across the transport. This design leans on the MCP contract to keep the LLM-facing surface stable while the internal computations evolve: adding a new loader or urgency signal does not change any tool signature. To our knowledge no other public MCP server exposes a persistent, cross-source work graph together with a bounded urgency queue and an inference-free persona through the same session; the closest analogues either publish a raw source (an IMAP MCP server) or a raw retriever (a vector-DB MCP server).

\subsection{Voice and style modelling}
\label{sec:related:voice}

Computational stylometry has a long history in author identification and forensic linguistics~\cite{TODO-stylometry-koppel,TODO-stamatatos-survey}, using features such as function-word frequencies, character n-grams, and average sentence length to fingerprint an author. On the generative side, recent work on style transfer and controllable generation~\cite{TODO-style-transfer-lample,TODO-controllable-generation} conditions an LLM on target-style examples or few-shot exemplars; commercial writing assistants (Grammarly, Lex) fine-tune on user text or use in-context exemplars to imitate a house voice. The intermediate approach --- extracting a compact statistical fingerprint and injecting it as a system-prompt bias --- is the one Karpathy's persona notes~\cite{TODO-karpathy-persona} popularised.

CoViber's \texttt{persona} module (\S\ref{sec:persona}) implements the intermediate approach in a form deliberately restricted to what closed-form statistics can produce. Given the corpus of records the user authored (i.e., records whose \texttt{from\_name} matches the configured identity), it computes seven scalar features --- message count, mean word count, formality score (a bounded ratio of formal to casual marker hits), emoji rate, question rate, exclamation rate --- plus the three most-common opener and closer bigrams and the twenty highest-frequency non-stopword, non-contraction signature words. These are wrapped into a system prompt of the form ``\emph{Write as $\langle$name$\rangle$. Voice: $\langle$register$\rangle$, $\sim$$k$ words per message, $\langle$emoji cue$\rangle$. Typically opens with `$X$' and closes with `$Y$'\ldots'', suitable for injection into any downstream drafting call. The design has three properties absent from LLM-based rewriters. First, it is \emph{deterministic}: identical input produces identical output, byte for byte, so a drafted reply can be regenerated and audited. Second, it is \emph{empty-corpus safe}: if no self-authored records are found, the module returns an explicit refusal-to-fabricate sentinel prompt rather than an implausible ``formal, $\sim$0 words per message'' description that would mislead callers. Third, it is \emph{cheap enough to run on every call}: no training step, no cached model, no GPU. The tradeoff is expressivity --- the profile cannot capture syntactic idiosyncrasies a fine-tuned rewriter could --- but the goal is stylistic guidance to a downstream LLM that already handles fluency, not full imitation.

\smallskip
\noindent In summary, CoViber's contribution is not any single component --- each of \S\ref{sec:related:rag}--\S\ref{sec:related:voice} has stronger prior art on one axis --- but the combination of a loader-agnostic core, a cross-source work model, inference-free scoring and voice modelling, and an MCP-native interface, packaged as a local-first, dependency-light artifact reproducible on synthetic data.



% -----------------------------------------------------------------------
%  §3  System
% -----------------------------------------------------------------------
% Verified: False
% Refutation notes:
%   The draft misstates the CLI surface. It writes: "The CLI (\texttt{coviber ingest / triage / recall / graph / voice / serve}, defined in \texttt{coviber/cli.py})". Checked against coviber/cli.py (build_parser, lines 121–145): the actual subcommands are `ingest`, `triage`, `query`, `graph`, `demo`, `loaders`, `serve`. There is no `recall` subcommand (that name is an MCP tool in mcp_server.py, not a CLI verb) and no `voice` subcommand (there is no CLI equivalent of the `voice_profile` MCP tool). The CLI also has `query`, `demo`, and `loaders` subcommands the draft omits. This is a straightforward factual error about a file the section explicitly references. Fix: either enumerate the real verbs (`coviber ingest / triage / query / graph / demo / loaders / serve`) or say something like "a set of one-shot verbs including \texttt{ingest}, \texttt{triage}, \texttt{query}, \texttt{graph}, and \texttt{serve}" without pretending `recall`/`voice` exist there. All other technical claims verified against source: pipeline.ingest matches Algorithm 1 line-for-line (get_loader→list→Store→upsert→all→WorkGraph(known_projects, you)→ingest→save_graph); 6 built-in loaders (demo/jsonl/mbox/imap/slackexport/webscrape) confirmed in loaders/__init__.py; `@register` decorator and `coviber.loaders` entry-point group confirmed; `Record.id`-keyed dedup with lock-guarded read-merge-rewrite and write-then-os.replace confirmed in store.upsert; `records.jsonl.bad` quarantine confirmed in _quarantine; four extras (`[search]`, `[scrape]`, `[qdrant]`, `[mcp]`) confirmed in pyproject.toml; JSONVectorStore vs QdrantVectorStore with `COVIBER_QDRANT_URL` env fallback confirmed in vector_stores.resolve; 7 MCP tools (recall, catch_me_up, who_is, project_status, graph_summary, voice_profile, refresh) all present as @mcp.tool() in mcp_server.py; per-call `_settings()` re-read confirmed in mcp_server._settings, and the "don't add a module-level settings cache" invariant is verbatim in ARCHITECTURE.md line 106; core stdlib-only claim holds (the only non-stdlib top-level core import is `mcp.server.fastmcp` in the [mcp]-gated mcp_server.py); "~10^5" JSON-store cutoff matches the vector_stores.py module docstring; inference-free persona confirmed (persona.py is pure statistics, no LLM calls). One minor caveat that does not itself warrant refutation: the draft's phrase "config-driven ticket regex" is technically true (WorkGraph.__init__ accepts `ticket_re=`), but Settings/pipeline.ingest do not currently thread a user-supplied regex through — only `known_projects` and `you` are threaded. Worth softening in prose but not a factual break.
% Notes to assembler: Architecture figure: the section references figures/architecture.pdf via \\includegraphics. Two options for the assembler:\n\n  (a) Render the ARCHITECTURE.md ASCII sketch as a TikZ diagram. I recommend a simple three-column layout: a left stack of six loader boxes (email/IMAP, mbox, slack export, j

\section{System}
\label{sec:system}

\subsection{Architecture overview}
\label{sec:arch}

CoViber is organized around a single architectural commitment: the only
pluggable seam is the boundary at which context enters the system. Everything
downstream --- deduplicated persistence, cross-source entity extraction,
urgency triage, semantic recall, and the client surface --- operates on a
common canonical schema and is deliberately source-agnostic. This section
gives the top-down view; the four numbered subsections that follow
(\S\ref{sec:records}--\S\ref{sec:persona}) formalize each stage in turn, and
\S\ref{sec:mcp} covers the client surface.

\begin{figure}[t]
\centering
\includegraphics[width=0.98\linewidth]{figures/architecture.pdf}
\caption{CoViber architecture. Heterogeneous sources are normalized by
per-platform \emph{Loaders} into a single \texttt{Record} stream. The
loader-agnostic core then applies four stages --- content-hash dedup and
persistence (\texttt{Store}), incremental entity extraction (\texttt{WorkGraph}),
multi-signal urgency triage, and inference-free voice modelling --- writing
three durable artefacts to disk: \texttt{records.jsonl}, a pluggable vector
index (\texttt{embeddings.json} or a Qdrant collection), and
\texttt{workgraph.json}. A CLI and an MCP stdio server share the same in-process
core and expose the memory to any MCP-capable client. All state is local; no
network egress is required for any core operation.}
\label{fig:architecture}
\end{figure}

\paragraph{The single seam.}
A CoViber \emph{Loader} is any object that yields \texttt{Record}s. The
\texttt{Record} dataclass (\S\ref{sec:records}) is the one shape every loader
produces; six built-ins ship in \texttt{coviber.loaders/} (\texttt{demo},
\texttt{jsonl}, \texttt{mbox}, \texttt{imap}, \texttt{slackexport},
\texttt{webscrape}), and third parties register additional loaders via a
\texttt{@register} decorator or the \texttt{coviber.loaders} entry-point
group.\footnote{See \texttt{coviber/loaders/base.py} and
\texttt{coviber/loaders/\_\_init\_\_.py}. The dispatch is name-keyed:
\texttt{get\_loader("myapp", cfg)} returns an instance regardless of whether the
loader lives inside the package or in a third-party wheel.} No downstream
module in \texttt{coviber/} branches on \texttt{record.source}; source-specific
parsing is confined to the loader module that produced the record.

\paragraph{Four core stages.}
A single call to \texttt{pipeline.ingest(settings)}
executes the loader-agnostic core as a linear sequence
(Algorithm~\ref{alg:ingest}):
\begin{enumerate}
    \item \textbf{Dedup and persistence.} \texttt{Store.upsert} merges the
    yielded records into \texttt{records.jsonl} keyed by the content-derived
    \texttt{Record.id} (\S\ref{sec:records}). Writers serialize on an advisory
    file lock (\texttt{.lock}) held across the read-merge-rewrite; readers stay
    lock-free because every rewrite is atomic (write-then-\texttt{os.replace}).
    Unparseable lines are quarantined to \texttt{records.jsonl.bad} rather than
    silently dropped~\cite{TODO-crash-consistency}.
    \item \textbf{Cross-source entity extraction.} \texttt{WorkGraph.ingest}
    walks the full record stream and folds people, projects, channels, and
    tickets into four bidirectionally-linked node tables using a
    config-driven ticket regex and case-insensitive canonicalization
    (\S\ref{sec:graph}). The result is persisted to \texttt{workgraph.json},
    again via a write-then-rename.
    \item \textbf{Urgency triage.} On demand --- typically from
    \texttt{pipeline.build\_queue} or the \texttt{catch\_me\_up} MCP tool ---
    the urgency scorer (\S\ref{sec:urgency}) applies a skip filter and a
    weighted multi-signal sum $U(r) \in [0,14]$ over the records loaded from
    disk, yielding a prioritized queue.
    \item \textbf{Semantic recall.} Also on demand, \texttt{Store.search}
    encodes a query and ranks records by cosine similarity against a
    persisted vector index; when the \texttt{[search]} extra is absent, it
    degrades to a keyword scorer over the same records
    (\S\ref{sec:recall}).
\end{enumerate}

\begin{algorithm}[t]
\caption{One ingest cycle (\texttt{coviber.pipeline.ingest}).}
\label{alg:ingest}
\begin{algorithmic}[1]
\Require settings $S$ (loader name, config, data dir, graph priors)
\State $L \gets \Call{GetLoader}{S.\text{loader}, S.\text{loader\_config}}$
\State $R \gets \Call{list}{L.\text{load}()}$ \Comment{loader-agnostic downstream}
\State $\Sigma \gets \Call{Store}{S.\text{data\_dir}, S.\text{qdrant}}$
\State $n_{\text{new}} \gets \Sigma.\Call{upsert}{R}$
    \Comment{lock-guarded read-merge-rewrite; dedup by \texttt{Record.id}}
\State $R_{\text{all}} \gets \Sigma.\Call{all}{}$ \Comment{single re-read; feeds both graph and stats}
\State $G \gets \Call{WorkGraph}{S.\text{known\_projects}, S.\text{you}}$
\State $G.\Call{ingest}{R_{\text{all}}}$ \Comment{$O(1)$ amortized per record}
\State $\Sigma.\Call{saveGraph}{G.\Call{toDict}{}}$
\State \Return $\{n_{\text{new}}, |R_{\text{all}}|, G.\Call{summary}{}\}$
\end{algorithmic}
\end{algorithm}

Note that triage and recall are \emph{not} part of the ingest cycle: they read
the persisted store and graph, so they can be invoked without a re-load. The
lock covers each writer's own read-merge-rewrite, not the whole ingest
pipeline --- two concurrent ingests on the same data directory interleave
safely at write boundaries but do not observe a transactional
``all-of-mine-then-all-of-yours'' view. This is a considered choice: a
whole-pipeline lock would push ingest latency up by the encoder cost, which
dominates on cold vectors.

\paragraph{Two persistence surfaces.}
The core writes to two logical stores, both under a single
\texttt{data\_dir}: (i)~\texttt{records.jsonl}, the append-optimized
canonical corpus of deduped \texttt{Record}s, plus its quarantine sibling
\texttt{records.jsonl.bad}; and (ii)~a pluggable vector index behind the
\texttt{VectorStore} protocol
(\texttt{coviber/vector\_stores.py}). Two backends implement the protocol:
\texttt{JSONVectorStore} materializes the whole vector table into one
\texttt{embeddings.json} file next to \texttt{records.jsonl} (zero
dependencies, appropriate up to $\sim 10^5$ records), and
\texttt{QdrantVectorStore} routes persistence and $k$-NN search to a Qdrant
instance (\texttt{[qdrant]} extra; typically a local Docker container).
Backend selection happens once at \texttt{Store} construction, keyed by an
explicit \texttt{qdrant:} block or the \texttt{COVIBER\_QDRANT\_URL} env
var; the search code path is identical regardless of backend. The
\texttt{workgraph.json} artefact is a third on-disk file but is fully derived
from \texttt{records.jsonl} and can be rebuilt at any time by re-running
\texttt{ingest}. See \S\ref{sec:recall} for the model-tag invalidation
protocol that keeps a persisted index consistent with the encoder in use.

\paragraph{Two client surfaces, one core.}
Both clients are thin wrappers over the same in-process modules
(Figure~\ref{fig:architecture}). The CLI (\texttt{coviber ingest / triage /
recall / graph / voice / serve}, defined in \texttt{coviber/cli.py}) is a
one-shot interface for humans and shell pipelines. The MCP server
(\texttt{coviber/mcp\_server.py}) is a long-running stdio process exposing
seven tools --- \texttt{recall}, \texttt{catch\_me\_up}, \texttt{who\_is},
\texttt{project\_status}, \texttt{graph\_summary}, \texttt{voice\_profile},
and \texttt{refresh} --- to any MCP-capable client
(\S\ref{sec:mcp})~\cite{TODO-mcp-spec}. Both routes go through
\texttt{Settings.from\_dict} and share a single YAML/JSON config schema
(\texttt{COVIBER\_CONFIG}), so an operator can migrate a session from
interactive CLI use to an always-on Claude-Desktop integration without
re-authoring configuration. The MCP server re-reads its settings on every
tool call, so a config edit takes effect on the next call without a
restart.\footnote{See the invariants note in
\texttt{ARCHITECTURE.md}: a module-level settings cache would break this
property and is explicitly disallowed.}

\paragraph{Design commitments.}
Four commitments shape the design and appear repeatedly in the sections that
follow. They are stated once here so later formalizations can refer to them
by name.
\begin{description}
    \item[Local-first.] Records, graph, and vectors live on disk under a
    user-owned \texttt{data\_dir}; the MCP transport is stdio to a local
    process. No network is required for any core operation. When Qdrant is
    configured, data still resides where the operator points the URL --- by
    default, a local Docker container.
    \item[One seam.] New sources are added by writing a new \texttt{Loader},
    not by changing core code. The invariant that ``no core module branches
    on \texttt{record.source}'' is the operational form of this commitment.
    \item[Graceful degradation.] The core has no non-stdlib runtime
    dependency. Four opt-in extras
    (\texttt{[search]}, \texttt{[scrape]}, \texttt{[qdrant]},
    \texttt{[mcp]}) each degrade cleanly when absent: embedding search
    falls back to keyword scoring, HTML scraping errors with an
    actionable message, the vector index defaults to
    \texttt{JSONVectorStore}, and the CLI still runs without
    \texttt{mcp}.
    \item[Inference-free voice.] The persona model
    (\S\ref{sec:persona}) is pure statistics --- no model call is made
    to draft in the user's style. Voice modelling is therefore fast,
    private, and deterministic on the local corpus.
\end{description}

The remainder of Section~\ref{sec:system} formalizes each stage. Sections
\S\ref{sec:records}--\S\ref{sec:recall} follow the pipeline order of
Algorithm~\ref{alg:ingest} (record schema, loader interface, work graph,
urgency triage, semantic recall), and \S\ref{sec:persona}--\S\ref{sec:mcp}
cover the persona model and the MCP client surface.

% Verified: False
% Refutation notes:
%   The draft misstates the warning category emitted on duplicate-name loader registration. Under "Registry", the draft claims: "if the same name is registered twice by different classes, the registry retains the newer binding and issues a `RuntimeWarning` identifying the shadowing module and class; silent override is disallowed." Verifying against `coviber/loaders/__init__.py` lines 18-21, the `register(name)` decorator's re-registration path invokes `warnings.warn(f"coviber loader '{name}' re-registered by ...", stacklevel=2)` with no `category=` argument, so the default `UserWarning` is emitted — not `RuntimeWarning`. This contradicts the code. (Note: the neighbouring claim about the *entry-point* failure warning being a `RuntimeWarning` is correct — line 70 explicitly passes `RuntimeWarning` — so only the duplicate-registration claim needs to be corrected, e.g. to `UserWarning` or the generic term "warning".) All other technical claims verified against code: 13-field Record ✓, id hash formula `MD5(text || \x1f || subject || \x1f || from_name || \x1f || source)` with `usedforsecurity=False` ✓ (record.py L50-54, L81), timestamp normalizer three-parser order with epoch sanity range [631152000, 4102444800] ✓ (record.py L30-42), naive→UTC and aware→astimezone(UTC) ✓ (L45-47), non-string coercion via `str()` ✓ (L27), `from_dict` unknown-key tolerance with per-key one-shot `RuntimeWarning` ✓ (L102-120), six built-in loaders demo/jsonl/mbox/imap/slackexport/webscrape ✓ (loaders/__init__.py L75), imap `BODY.PEEK` + read-only + timeout + `password_env` + plaintext-password rejection ✓ (imap.py L30-77), mbox+imap positive-limit rejection ✓, slack export users.json + `<@Uxxx>` mention resolution + replied-from-thread-membership ✓ (slackexport.py L27,L80-95), webscrape `record_selector` + `fields{selector, attr}` + `bs4` behind `[scrape]` extra ✓ (webscrape.py L56-58). One additional minor imprecision worth flagging (not a refutation on its own but tied to the fix): the draft's tuple/table uses the mathematical short names `from` and `thread`, while the actual dataclass fields are `from_name` and `thread_id` (record.py L68, L74). The draft frames these as a formal schema so the abbreviation is defensible, but an assembler note would help readers cross-referencing the code.
% Notes to assembler: Prompt vs. code alignment notes:\n\n1. The prompt says 'EXCLUDE channel and ts from the id'. The code excludes considerably more than that: only (text, subject, from_name, source) participate in the hash (see record.py line 81). I stated the full exclusion list (channel, ts, unread, replied, recipie

\subsection{Record schema}
\label{sec:record}

Every source of context---email, chat, code review, issue tracker, or arbitrary
web page---is projected into a single canonical structure, the
\emph{Record}. Downstream stages (deduplication, graph, urgency, search,
persona) read Records and only Records; source-specific parsing lives strictly
upstream of this boundary (\S\ref{sec:loader}). Formally, a Record is a tuple
\[
r \;=\; (\mathit{source},\, \mathit{text},\, \mathit{subject},\, \mathit{from},\, \mathit{recipient},\, \mathit{channel},\, \mathit{url},\, \mathit{ts},\, \mathit{unread},\, \mathit{thread},\, \mathit{replied},\, \mathit{scraped\_at},\, \mathit{id})
\]
over the domains summarised in Table~\ref{tab:record}. Only $\mathit{source}$
is required in practice; all other fields default and are populated
opportunistically by whichever loader produced the record.

\begin{table}[h]
\centering
\small
\begin{tabular}{@{}lll@{}}
\toprule
Field & Type & Role \\
\midrule
$\mathit{source}$      & string  & loader identifier, e.g.\ \texttt{email}, \texttt{slack}, \texttt{github} \\
$\mathit{text}$        & string  & body / message content \\
$\mathit{subject}$     & string  & title / subject line \\
$\mathit{from}$        & string  & sender display name \\
$\mathit{recipient}$   & string  & primary addressee (typically \texttt{"you"}) \\
$\mathit{channel}$     & string  & sub-view (\texttt{dm}, \texttt{inbox}, \texttt{\#general}, \dots) \\
$\mathit{url}$         & string  & deep link, if the source exposes one \\
$\mathit{ts}$          & string  & event timestamp, ISO-8601 UTC (see below) \\
$\mathit{unread}$      & bool    & has the user viewed this item \\
$\mathit{thread}$      & string  & thread identifier, for reply detection \\
$\mathit{replied}$     & bool    & has the user replied in this thread \\
$\mathit{scraped\_at}$ & string  & wall-clock of ingest \\
$\mathit{id}$          & string  & content-derived deduplication key \\
\bottomrule
\end{tabular}
\caption{The thirteen fields of a canonical Record.}
\label{tab:record}
\end{table}

\paragraph{Identity by content.}
The record identifier is derived deterministically from a fixed prefix of the
tuple. Let $\mathrm{H}$ be MD5 and let $\us$ denote the ASCII unit separator
\texttt{\textbackslash x1f}. Then
\begin{equation}
\mathit{id}(r) \;=\; \mathrm{H}\bigl(\, \mathit{text}(r) \,\Vert\, \us \,\Vert\, \mathit{subject}(r) \,\Vert\, \us \,\Vert\, \mathit{from}(r) \,\Vert\, \us \,\Vert\, \mathit{source}(r) \,\bigr).
\label{eq:id}
\end{equation}
Unset fields participate as the empty string. Two properties of
Eq.~(\ref{eq:id}) matter downstream.
\emph{(i)} The $\us$ delimiter guarantees an unambiguous field boundary:
$(\text{``ab''},\text{``c''})$ and $(\text{``a''},\text{``bc''})$ hash to
different ids, closing the class of concatenation collisions that a naive
join would admit.
\emph{(ii)} MD5 is invoked with \texttt{usedforsecurity=False}: the digest is
used only as a dedup key, never as a cryptographic commitment, and the flag
lets the same code run under FIPS-mode Python. Any collision-resistant hash
would work; MD5 was retained for its narrow (32-hex-char) representation and
zero-dependency availability.

\paragraph{Dedup invariant.}
By construction, Eq.~(\ref{eq:id}) yields the
\emph{Record dedup invariant}: two records $r_1, r_2$ with identical
$(\mathit{text}, \mathit{subject}, \mathit{from}, \mathit{source})$ satisfy
$\mathit{id}(r_1) = \mathit{id}(r_2)$. The store keys on $\mathit{id}$
(\S\ref{sec:store}), so re-scraping the same source repeatedly is idempotent:
the corpus converges rather than growing linearly in scrape count. This is the
single property that makes ``re-run the loaders whenever'' a viable
operational mode.

\paragraph{Fields deliberately excluded from the id.}
The hash covers content and origin only; it omits $\mathit{channel}$,
$\mathit{ts}$, $\mathit{unread}$, $\mathit{replied}$, $\mathit{recipient}$,
$\mathit{url}$, $\mathit{thread}$, and $\mathit{scraped\_at}$. Two exclusions
are load-bearing and documented as
architectural invariants~\cite{TODO-coviber-architecture-md}:
\begin{itemize}
  \item \textbf{$\mathit{channel}$ excluded.} The same message surfaced through
    two views (e.g.\ a Slack DM cross-posted to a channel, or an email visible
    in both \texttt{inbox} and a label) collapses to one record. Escaping
    this collapse when it is unwanted is a per-loader concern, addressed by
    varying $\mathit{source}$ (issue tracker~\cite{TODO-coviber-issue-17}).
  \item \textbf{$\mathit{ts}$ and mutable state excluded.} A CI bot that emits
    an identical ``build passed'' line daily produces a single record whose
    $\mathit{ts}$ advances on re-scrape; likewise $\mathit{replied}$ and
    $\mathit{unread}$ can flip without minting a new id. The tradeoff is that
    genuinely-recurring identical events are undercounted; the benefit is
    that mutable state converges to its latest observation rather than
    forking into a new record on every state transition.
\end{itemize}

\paragraph{Timestamp normalisation.}
$\mathit{ts}$ is normalised in \texttt{Record.\_\_post\_init\_\_} so that
lexicographic ordering of the stored string equals chronological ordering.
The normalisation function attempts three parsers in order:
\emph{(1)} ISO-8601 (with a \texttt{Z}$\rightarrow$\texttt{+00:00} shim, since
Python~3.9's \texttt{datetime.fromisoformat} does not accept the military
suffix); \emph{(2)} RFC-2822 (the format email \texttt{Date:} headers use),
via \texttt{email.utils.parsedate\_to\_datetime}; \emph{(3)} POSIX epoch
seconds, gated to the sanity range $[631152000, 4102444800]$ (1990-01-01 to
2100-01-01) so that small integers are not misinterpreted as timestamps.
Naive datetimes are assumed UTC; aware datetimes are converted via
\texttt{astimezone(UTC)} so mixed-offset inputs still order correctly.
Non-string inputs (int/float epoch, bool, \texttt{None}) are coerced through
\texttt{str()} before parsing---a defensive choice motivated by JSONL inputs
where the surface syntax admits any of these. Unparseable strings pass through
unchanged; the loss is confined to that record's ordering, not the store's
readability.

\paragraph{Forward-compatible deserialisation.}
\texttt{Record.from\_dict} tolerates unknown keys: an older CoViber reading a
\texttt{records.jsonl} written by a newer version silently drops fields it
cannot represent, rather than crashing. The first observation of each unknown
key raises a per-process \texttt{RuntimeWarning}, so an operator running a
downgraded binary is told which fields will be lost on the next rewrite. Warn
frequency is capped at once per key to prevent a large store from drowning
real signals in duplicate warnings.

\subsection{Loader interface}
\label{sec:loader}

The Loader is the \emph{only} extension seam in CoViber: everything downstream
of the Record boundary---dedup, graph, urgency, persona, MCP---is
source-independent by construction. Adding a data source therefore reduces to
implementing a single method.

\paragraph{Contract.}
A Loader is a Python class satisfying
\begin{equation*}
\texttt{Loader.load()} \;:\; \varnothing \;\longrightarrow\; \mathrm{Iterable}[\mathrm{Record}].
\end{equation*}
The default constructor accepts a \texttt{dict} of configuration (the block
under \texttt{config:} in the YAML file, or a keyword equivalent). Loaders are
expected to validate required configuration eagerly in \texttt{\_\_init\_\_}
and to stream records lazily from \texttt{load()}, so that a large corpus
(a full mbox, an IMAP inbox, a long JSONL file) can be ingested without
holding it in memory. The base class provides no default implementation of
\texttt{load()}; the class attribute \texttt{name} is populated by the
registry decorator (below), giving a loader a stable identifier that YAML
config and the CLI can reference.

\paragraph{Registry.}
Loaders resolve by name through a process-global registry. The
\texttt{@register("myapp")} decorator inserts a class into the registry under
that name and sets its \texttt{name} attribute. Third-party packages may
alternatively expose loaders via the standard Python \texttt{coviber.loaders}
entry-point group, which is loaded lazily the first time the registry is
queried. Resolution proceeds as follows:

\begin{algorithm}[h]
\caption{\textsc{get\_loader}$(n, c)$ --- resolve a Loader by name}
\label{alg:get_loader}
\begin{algorithmic}[1]
\Require name $n$, config dict $c$
\State \Call{load\_entrypoints}{} \Comment{once per process; discovers third-party loaders}
\If{$n \notin \mathrm{REGISTRY}$}
  \State \textbf{raise} $\mathrm{KeyError}(n, \mathrm{available}=\mathrm{sorted}(\mathrm{REGISTRY}))$
\EndIf
\State \Return $\mathrm{REGISTRY}[n](c)$
\end{algorithmic}
\end{algorithm}

Two properties matter. First, if the same name is registered twice by
different classes, the registry retains the newer binding and issues a
\texttt{RuntimeWarning} identifying the shadowing module and class; silent
override is disallowed. Second, entry-point loading is \emph{fail-open}: a
broken third-party plugin (import error, missing extra) does not disable the
whole registry, but does emit a \texttt{RuntimeWarning} naming the entry point
and the underlying exception. Silent failure would be catastrophic for an
extension seam whose value proposition is ``write your own loader in ${\sim}10$
lines''; the operator must be able to see \emph{why} a plugin did not
register.

\paragraph{Built-in loaders.}
Six loaders ship in the reference implementation
(Table~\ref{tab:builtin-loaders}); each is under $\sim\!100$ lines of
pure-standard-library Python, with the sole exception of
\texttt{webscrape} which imports \texttt{beautifulsoup4} behind the
\texttt{[scrape]} extra.

\begin{table}[h]
\centering
\small
\begin{tabular}{@{}lp{0.72\linewidth}@{}}
\toprule
Name & Description \\
\midrule
\texttt{demo}        & Fixed synthetic cross-platform stream for a fictional company; zero configuration. Powers the reproducible evaluation of \S\ref{sec:eval}. \\
\texttt{jsonl}       & Universal escape hatch: one JSON object per line, keys are a subset of the Record fields. Any data one can dump to disk can enter CoViber through this loader. \\
\texttt{mbox}        & Local Unix \texttt{mbox} file via \texttt{mailbox} (stdlib). Newest-first, with a positive-integer \texttt{limit} for the newest $N$ messages. Ships a stdlib HTML$\rightarrow$text fallback for HTML-only messages. \\
\texttt{imap}        & IMAP4-over-SSL fetch. Password is referenced by environment-variable \emph{name} (\texttt{password\_env}); plaintext passwords in config are rejected. Mailbox is opened read-only; \texttt{BODY.PEEK} avoids marking mail as seen. A configurable network timeout bounds every syscall. \\
\texttt{slackexport} & A standard Slack workspace export (unzipped directory of \texttt{users.json} plus per-channel day files). Resolves \texttt{<@Uxxx>} mentions to display names; derives \texttt{replied} from thread membership. \\
\texttt{webscrape}   & Config-driven DOM extraction: a \texttt{record\_selector} plus a map of field$\rightarrow$(CSS selector, attribute) turns any list-style web page into Records with zero per-site code. \\
\bottomrule
\end{tabular}
\caption{Built-in loaders. Every non-trivial parsing seam CoViber needs is one of these six; a seventh is a two-file plugin.}
\label{tab:builtin-loaders}
\end{table}

Two structural choices deserve emphasis. The \texttt{jsonl} loader is
sufficient by itself to onboard any source that can be serialised to disk;
this is the ``no plugin required'' path for exotic corpora. The
\texttt{webscrape} loader---where the site's DOM shape lives in YAML rather
than Python---is the ``no plugin required'' path for browser-visible content.
Together they cover the long tail without expanding the registry.

\paragraph{Consequence for the rest of the system.}
Because the Loader is the sole seam and every Loader emits Records satisfying
Eq.~(\ref{eq:id}), the deduplication invariant, the WorkGraph
(\S\ref{sec:workgraph}), the urgency scoring function
(\S\ref{sec:urgency}), the persona model (\S\ref{sec:persona}), and the
MCP surface (\S\ref{sec:mcp}) each admit a single implementation valid
across all present and future sources. This is the load-bearing simplification
behind CoViber's small line count: the source-heterogeneity problem is
contained entirely in \S\ref{sec:loader}.

% Verified: True
% Notes to assembler: Preamble deps: algorithm/algorithmicx (or algpseudocode), booktabs (for the schema table), siunitx (for the placeholder), amsmath, amssymb. No figure files needed. Cross-refs used and expected to resolve elsewhere: sec:vector-store, sec:record, sec:urgency, sec:pipeline, sec:eval, sec:audit, sec:mcp

\section{The Cross-Source WorkGraph}
\label{sec:workgraph}

The WorkGraph is CoViber's structural memory: a heterogeneous entity graph induced from the record stream that unifies people, projects, channels, and tickets scattered across loaders. Where the vector index (\S\ref{sec:vector-store}) supports lexical and semantic recall over individual records, the WorkGraph answers relational queries — \emph{who} works with whom, \emph{which} projects a channel touches, \emph{which} tickets a person has been active on — that no per-record embedding can serve directly. It is populated purely from the fields of the canonical \texttt{Record} schema (\S\ref{sec:record}); no loader-specific logic reaches into \texttt{workgraph.py}.

\subsection{Definitions}
\label{sec:workgraph-defs}

Let $\mathcal{R}$ be the deduplicated record store and let $\mathcal{P}_{\text{known}}$ be an operator-supplied set of project names (config-driven; empty is legal, in which case the projects partition is empty). Let $Y$ be the operator's own identity (the \texttt{you} setting), matched case-insensitively.

The WorkGraph is a tuple $G = (V_{\mathrm{peo}}, V_{\mathrm{prj}}, V_{\mathrm{ch}}, V_{\mathrm{tk}}, E)$ where the four node partitions are:
\begin{itemize}
  \setlength\itemsep{0.15em}
  \item $V_{\mathrm{peo}}$: distinct persons observed across all records, keyed by their case-folded name.
  \item $V_{\mathrm{prj}}$: elements of $\mathcal{P}_{\text{known}}$ that appeared in at least one record.
  \item $V_{\mathrm{ch}}$: distinct non-empty values of the \texttt{channel} field.
  \item $V_{\mathrm{tk}}$: distinct ticket / PR / issue identifiers extracted from record text and URL.
\end{itemize}
Edges are induced by co-occurrence within a record. Formally, for each record $r \in \mathcal{R}$ let $\mathrm{peo}(r), \mathrm{prj}(r), \mathrm{ch}(r), \mathrm{tk}(r)$ denote the entities extracted from $r$ (defined operationally in \S\ref{sec:workgraph-alg}). Then $E$ contains an undirected edge between every pair drawn from $\mathrm{peo}(r) \cup \mathrm{prj}(r) \cup \mathrm{ch}(r) \cup \mathrm{tk}(r)$ that crosses partitions, except that no edges are induced \emph{within} the tickets partition or \emph{within} the channels partition (mirroring the adjacency-set layout in \texttt{workgraph.py}). Each node additionally carries scalar attributes: an interaction count (people only), a mention count (projects, channels, tickets), a \texttt{last\_seen} timestamp (people only), a \texttt{display\_name} (people only), and the set of platforms it has appeared on (people only). The precise adjacency-set schema per partition is:
\begin{center}
\begin{tabular}{lll}
\toprule
partition & adjacency sets & scalars \\
\midrule
people    & platforms, projects, channels & interaction\_count, last\_seen, display\_name \\
projects  & people, channels, tickets     & mentions \\
channels  & people, projects              & mentions \\
tickets   & people, projects              & mentions \\
\bottomrule
\end{tabular}
\end{center}

The graph is intentionally undirected. Recipient / sender asymmetry, thread role, and reply direction are already carried on the underlying \texttt{Record} and consumed by the urgency stage (\S\ref{sec:urgency}); duplicating them as edge attributes here would double the storage cost of the graph for a signal the triage layer already owns.

\subsection{Entity extraction}
\label{sec:workgraph-alg}

Algorithm~\ref{alg:workgraph-ingest} gives the per-record extraction and mutation step, applied by \texttt{WorkGraph.\_ingest\_one} to each $r \in \mathcal{R}$ in a single pass.

\begin{algorithm}[t]
\caption{WorkGraph ingest: single record}
\label{alg:workgraph-ingest}
\begin{algorithmic}[1]
\Require record $r$; known projects $\mathcal{P}_{\text{known}}$ (pre-compiled to word-boundary patterns); ticket regex $T$; identity $Y$
\Ensure in-place mutation of $V_{\mathrm{peo}}, V_{\mathrm{prj}}, V_{\mathrm{ch}}, V_{\mathrm{tk}}$
\State $\mathit{blob} \gets r.\mathrm{subject} \Vert \text{`` ''} \Vert r.\mathrm{text}$
\State $\mathit{blob}_\ell \gets \mathrm{lower}(\mathit{blob})$;\ \ $y \gets \mathrm{lower}(Y)$
\Statex \Comment{Person extraction}
\State $S \gets \{r.\mathrm{from\_name}, r.\mathrm{recipient}\} \setminus \{\varepsilon\}$ \Comment{drop empty strings}
\State $S \gets S \cup \{m \mid m \in \textsc{MentionRe}.\mathrm{findall}(r.\mathrm{text})\}$
\State $\mathit{peo} \gets \emptyset$; \ $\mathit{display} \gets \{\}$ \Comment{lowercase key $\to$ best display form}
\ForAll{$p \in S$}
  \State $k \gets \mathrm{lower}(p)$
  \If{$k = y$} \textbf{continue} \EndIf
  \If{$k \notin \mathit{display}$ \textbf{or} ($\mathit{display}[k] = k$ \textbf{and} $p \neq k$)}
    \State $\mathit{display}[k] \gets p$
  \EndIf
  \State $\mathit{peo} \gets \mathit{peo} \cup \{k\}$
\EndFor
\Statex \Comment{Project matching (pre-compiled, word-boundary)}
\State $\mathit{prj} \gets \{\, p : (p, \pi) \in \Pi \ \land\ \pi.\mathrm{search}(\mathit{blob}_\ell) \,\}$
\Statex \Comment{Channel}
\State $\mathit{ch} \gets \{r.\mathrm{channel}\}$ if $r.\mathrm{channel} \neq \varepsilon$ else $\emptyset$
\Statex \Comment{Ticket extraction and canonicalization}
\State $M \gets T.\mathrm{findall}(\mathit{blob} \Vert \text{`` ''} \Vert r.\mathrm{url})$
\State $\mathit{tk} \gets \{\, \mathrm{sub}(\mathtt{r"PR\textbackslash s*\#\textbackslash s*"},\ \mathtt{"PR\ \#"},\ m)\ :\ m \in M\, \}$
\Statex \Comment{Node mutation}
\ForAll{$k \in \mathit{peo}$}
  \State update $V_{\mathrm{peo}}[k]$: refresh \texttt{display\_name} if new form is better; $\mathrm{platforms} \gets \mathrm{platforms} \cup \{r.\mathrm{source}\}$;
  \Statex\hspace{2em}$\mathrm{projects} \gets \mathrm{projects} \cup \mathit{prj}$;\ \ $\mathrm{channels} \gets \mathrm{channels} \cup \mathit{ch}$;
  \Statex\hspace{2em}$\mathrm{interaction\_count} \mathrel{+}= 1$;\ \ update $\mathrm{last\_seen}$ if $r.\mathrm{ts}$ is ISO
\EndFor
\ForAll{$p \in \mathit{prj}$}\ update $V_{\mathrm{prj}}[p]$: $\mathrm{people} \cup \mathit{peo}$, $\mathrm{channels} \cup \mathit{ch}$, $\mathrm{tickets} \cup \mathit{tk}$, $\mathrm{mentions}\mathrel{+}=1$ \EndFor
\ForAll{$c \in \mathit{ch}$}\ update $V_{\mathrm{ch}}[c]$: $\mathrm{people} \cup \mathit{peo}$, $\mathrm{projects} \cup \mathit{prj}$, $\mathrm{mentions}\mathrel{+}=1$ \EndFor
\ForAll{$t \in \mathit{tk}$}\ update $V_{\mathrm{tk}}[t]$: $\mathrm{people} \cup \mathit{peo}$, $\mathrm{projects} \cup \mathit{prj}$, $\mathrm{mentions}\mathrel{+}=1$ \EndFor
\end{algorithmic}
\end{algorithm}

Four extractors deserve individual comment because each fixes a class of failure the naive implementation exhibited in earlier revisions of the system.

\paragraph{Person extraction (lines 3--12).}
The candidate set is the union of the record's \texttt{from\_name}, its \texttt{recipient}, and every \texttt{@}-handle extracted from the body by the mention regex \texttt{MentionRe}:
\[
\mathtt{(?<![\backslash w.])@([A\text{-}Za\text{-}z0\text{-}9\_](?:[A\text{-}Za\text{-}z0\text{-}9\_.\text{-}]*[A\text{-}Za\text{-}z0\text{-}9\_])?)}
\]
The negative lookbehind \texttt{(?<![\textbackslash w.])} rejects the local-part-plus-\texttt{@} prefix of an email address (so \texttt{alice@example.com} does not spuriously mint an entity called \emph{example}), and the trailing character class forbids a handle ending in \texttt{.} or \texttt{-}, so sentence-final \texttt{@ada.} yields \texttt{ada} rather than \texttt{ada.}.

Person node keys are case-folded so \texttt{"Ada Byron"} (from an email header) and \texttt{"ada byron"} (from a Slack export) collapse to one node. The pre-fold display form is retained on the node as \texttt{display\_name} under a first-good-form-wins rule: a mixed-case form replaces a lowercase-only placeholder, but a subsequent all-uppercase form does not overwrite an earlier mixed-case one. The operator's own identity $Y$ is filtered out at extraction time to prevent the graph collapsing into a hub-and-spoke topology centred on the user (every co-occurrence edge would otherwise touch $Y$). The current implementation does \emph{not} attempt to reconcile \texttt{"Ada Byron"} with the mention handle \texttt{@ada}; alias resolution would require an operator-supplied handle table and is deferred to future work (\S\ref{sec:future-work}).

\paragraph{Project matching (line 13).}
For each $p \in \mathcal{P}_{\text{known}}$ the constructor compiles the pattern $\pi_p = \mathtt{re.compile(rf"\textbackslash b\{re.escape(p.lower())\}\textbackslash b")}$ once, and $r$ is matched against the case-folded blob. Two design choices are worth calling out.

First, \emph{word-boundary rather than substring}. An earlier revision used a plain \texttt{in} test on the lowercase blob, which caused short project names to produce large numbers of false-positive edges: \texttt{"AI"} matched \emph{email}, \emph{rail}, \emph{tail}; \texttt{"Go"} matched \emph{going}, \emph{ago}, \emph{cargo}. Word-boundary anchoring resolves this at the cost of missing intentional CamelCase compounds like \texttt{"CoViberAI"}, which we judge acceptable: operators who need such matches can simply add the compound to $\mathcal{P}_{\text{known}}$.

Second, \emph{pre-compilation}. The alternative — recompiling the $|\mathcal{P}_{\text{known}}|$ patterns per record — turned regex compilation into a first-order cost at $|\mathcal{R}| = 10^5$ in early profiling. \texttt{re.escape} ensures that operational project names containing regex metacharacters (\texttt{"C++"}, \texttt{".NET"}, \texttt{"F\#"}) are treated as literals rather than as patterns.

\paragraph{Ticket extraction (lines 14--15).}
The default ticket pattern is
\[
\mathtt{(\backslash b[A\text{-}Z][A\text{-}Z0\text{-}9]+\text{-}\backslash d+\backslash b\ |\ \backslash bPR\backslash s*\#\backslash d+\backslash b\ |\ (?<![\backslash w\#])\#\backslash d\{3,\}\backslash b)}
\]
which admits three families: uppercase project-prefixed tickets (\texttt{"PROJ-123"}, \texttt{"AB1-42"}), pull-request references with optional whitespace between \texttt{"PR"} and \texttt{"\#"} (\texttt{"PR \#12"}, \texttt{"PR\#12"}), and bare-hash issue numbers with at least three digits (\texttt{"\#4567"}). The 3-digit minimum on the bare-hash branch is deliberate: it filters out ordinal fragments (\texttt{"\#1"}, \texttt{"\#42"}) that appear in list markup and casual prose without meaning a ticket, at the cost of missing short issue IDs early in a project. The lookbehind \texttt{(?<![\textbackslash w\#])} rejects both word-internal hashes and doubled hashes such as those found in Markdown H2 headers. The pattern is applied to \texttt{blob $\Vert$ url}, so PR references that appear only as a link in \texttt{r.url} are still captured.

Every match is passed through \texttt{re.sub(r"PR\textbackslash s*\#\textbackslash s*", "PR \#", m)}, canonicalizing the whitespace-permissive branch to a single form. This prevents \texttt{"PR \#482"}, \texttt{"PR\#482"}, and \texttt{"PR  \#482"} from becoming three distinct nodes when the same PR is re-scraped through different loaders — a duplication bug the canonicalization step exists to close.

\paragraph{Channels (line 16).}
Channels are a passthrough of the \texttt{Record.channel} field. Loaders normalize their platform-native container concept into this field (a Slack channel name, a mail folder, a GitHub repository, a chat DM identifier), so the WorkGraph does not need to know which loader produced the record.

\subsection{Complexity and current re-build model}
\label{sec:workgraph-complexity}

Per-record extraction is $O(1)$ in $|\mathcal{R}|$: the person and mention regexes scan a bounded record; project matching scans the blob once per known project ($O(|\mathcal{P}_{\text{known}}|)$, with $|\mathcal{P}_{\text{known}}|$ typically $\leq 50$); ticket extraction is one linear scan; set updates are amortized $O(1)$ against Python's hash-set primitives. Full-corpus ingest is therefore $\Theta(|\mathcal{R}|)$ with a small constant.

The current pipeline (\texttt{pipeline.ingest}, \S\ref{sec:pipeline}) rebuilds the graph from scratch on every ingest cycle: after \texttt{Store.upsert}, it reads the full deduplicated corpus back with \texttt{store.all()} and calls \texttt{WorkGraph.ingest(all\_records)}. This is deliberate under the current design point, where $|\mathcal{R}| = 10^4$--$10^5$ and a full rebuild completes in $\approx\SI{XXX}{\milli\second}$ on the reference machine (\S\ref{sec:eval}); it also elides the correctness question of how to \emph{remove} edges when a record is superseded by a re-scrape with corrected metadata, since re-scrapes never mutate old edges in a rebuild.

For target deployments in the $|\mathcal{R}| > 10^6$ regime this becomes the dominant per-cycle cost. An incremental variant that (i) processes only the new records returned by \texttt{Store.upsert} and (ii) reference-counts edges so re-scrapes can decrement is tracked as issue \cite{TODO-issue-18} and briefly discussed in \S\ref{sec:future-work}.

\subsection{Serialization and query surface}
\label{sec:workgraph-queries}

\texttt{WorkGraph.to\_dict} materializes the four partitions as JSON-friendly nested dictionaries: each adjacency \emph{set} is sorted for deterministic output (so the persisted graph diffs cleanly across ingest runs, which matters for the audit-log design of \S\ref{sec:audit}). Two query helpers are exposed directly on the class:
\begin{itemize}
  \setlength\itemsep{0.15em}
  \item \texttt{person(name)} and \texttt{project(name)} return the serialized node for a single entity, used by the MCP tools \texttt{who\_is} and \texttt{project\_status} (\S\ref{sec:mcp}).
  \item \texttt{summary()} returns partition sizes, the top-10 people by interaction count, and the sorted project list; this backs the \texttt{graph\_summary} MCP tool and is what \texttt{pipeline.ingest} returns to its caller for status reporting.
\end{itemize}
Multi-hop queries (``which projects does person $X$ share with person $Y$?'', ``which tickets on project $P$ has person $X$ touched?'') are answered by set intersection over the adjacency fields, again in $O(1)$ per adjacency lookup and $O(|V|)$ in the worst case; no formal graph traversal library is used or needed at the current scale.


% Verified: True
% Notes to assembler: Preamble deps needed: \\usepackage{algorithm}, \\usepackage{algpseudocode}, \\usepackage{booktabs}, \\usepackage{amsthm} with a `definition` theorem style (Definition 4 = urgency, Definition 5 = skip filter — numbering is intentional per the prompt), \\usepackage{siunitx} (not used here but consiste

\subsection{Multi-Signal Urgency Triage}
\label{sec:urgency}

Once the store is populated and enriched with a \WorkGraph{} (\S\ref{sec:workgraph}),
the remaining problem is inbox-scale: how does the system decide what,
among $N$ heterogeneous records, actually deserves the user's attention
this hour?  \coviber{}'s answer is a rule-based additive scorer with a
prefilter, implemented in \path{coviber/urgency.py}. The design goal is
explicit: the scoring function must be inspectable, deterministic, and
reproducible on-device, with no reliance on a hosted classifier and no
labeled training corpus.

\subsubsection{Formal definition}

Let $\mathcal{S} = \{s_1, \ldots, s_{10}\}$ be the fixed set of signal
predicates enumerated in Table~\ref{tab:signals}. Each signal is a
Boolean predicate $s_i : \mathcal{R} \to \{0,1\}$ over a record $r$
(cf.\ Definition~\ref{def:record}). Let $w : \mathcal{S} \to
\mathbb{Z}_{\ge 0}$ be a configured weight map. The three age-tier
signals $\{s_{\text{age\_7d}}, s_{\text{age\_48h}}, s_{\text{age\_24h}}\}$
are mutually exclusive by construction (they are the branches of a
single \texttt{if/elif/elif} on record age), so at most one fires per
record.

\begin{definition}[Urgency score]
\label{def:urgency}
The urgency of a record $r$ is
\[
U(r) \;=\; \sum_{s \in S(r)} w(s), \qquad
S(r) \;=\; \{\, s \in \mathcal{S} \mid s(r) = 1 \,\},
\]
with a signal contributing zero (and being treated as disabled) iff
$w(s) \le 0$.
\end{definition}

At the default weights shipped with the package
(Table~\ref{tab:signals}), the seven non-age signals sum to $11$ and the
maximal age tier contributes $3$, so $U(r) \in [0, 14]$.  This
$[0,14]$ contract is pinned by a regression test
(\texttt{test\_urgency\_contract}); no clamp is applied at runtime.
Under a custom weight map the ceiling shifts in proportion to the
caller's choices, which is deliberate: silent clamping would mask a
misconfiguration rather than surface it. Configs that supply only a
partial weight dictionary are merged key-by-key against the defaults,
unknown keys are surfaced as \texttt{RuntimeWarning}, and non-integer
values raise a \texttt{ValueError} that names the offending key. Weight
$0$ is a first-class way to disable a signal.

\begin{definition}[Skip filter]
\label{def:skip}
A record $r$ is \emph{skipped} (excluded from the queue before scoring)
iff
\[
\mathrm{skip}(r) \;=\; \mathrm{skip\_sender}(r) \;\lor\; \mathrm{skip\_subject}(r) \;\lor\; \mathrm{fyi}(r),
\]
where $\mathrm{skip\_sender}$ matches any element of
$\mathtt{cfg.skip\_senders}$ against \texttt{r.from\_name} on token
boundaries, $\mathrm{skip\_subject}$ matches any element of
$\mathtt{cfg.skip\_subjects}$ as a case-insensitive substring of
\texttt{r.subject}, and $\mathrm{fyi}$ matches any element of the FYI
lexicon (defaults: \texttt{fyi}, \texttt{no action needed},
\texttt{no reply needed}, \texttt{just so you know}) against the
concatenation of subject and body, again on token boundaries.
\end{definition}

The token-boundary discipline for both $\mathrm{skip\_sender}$ and
$\mathrm{fyi}$ is load-bearing. Prior versions of \coviber{} matched
these lexicons as raw substrings, which caused two classes of false
positive: (i) colleagues whose surname literally contained a skipped
token (e.g. \texttt{"Abbott"} matching \texttt{"bot"}), and (ii)
common English words containing the trigram \texttt{"fyi"} inside them
(\texttt{"justifying"}, \texttt{"notifying"}, \texttt{"typify"}). Both
would silently drop legitimate obligations from the queue. The v0.4/L5
audit tightened both matches to word-boundary regexes of the form
\verb|(?<![a-z0-9])<lex>(?![a-z0-9])|; the fix is small, but it is the
canonical example of the transparency guarantee that motivates
rule-based scoring in the first place.

The complete triage procedure is given in Algorithm~\ref{alg:triage}.
Three additional filters run in the loop but are not part of
Definition~\ref{def:skip}: self-authored records (where
\texttt{r.from\_name} equals \texttt{cfg.you}) are dropped as
non-obligations, the skip filter runs first, and any record that scores
$U(r) \le 0$ after scoring is dropped as noise. The surviving records
are sorted by $(-U(r), r.\mathrm{ts})$ so that ties break in favour of
older records.

\begin{algorithm}
\caption{Triage: filter, score, sort.}
\label{alg:triage}
\begin{algorithmic}[1]
\Require records $R$, urgency config $\mathrm{cfg}$
\Ensure ranked queue $Q$ of $\{\mathrm{record}, U, S(r)\}$
\State $Q \gets [\,]$
\ForAll{$r \in R$}
  \If{$r.\mathrm{from\_name}.\mathrm{lower}() = \mathrm{cfg.you}.\mathrm{lower}()$} \State \textbf{continue} \EndIf
  \If{$\mathrm{skip}(r, \mathrm{cfg}) \neq \varnothing$} \State \textbf{continue} \EndIf
  \State $(u, S) \gets \mathrm{score}(r, \mathrm{cfg})$
  \If{$u \le 0$} \State \textbf{continue} \EndIf
  \State $Q.\mathrm{append}(\{\mathrm{record}\!:\!r,\, U\!:\!u,\, S\!:\!S\})$
\EndFor
\State \Return $Q$ sorted by $(-U, r.\mathrm{ts})$
\end{algorithmic}
\end{algorithm}

Complexity is $O(|R| \cdot k)$ per call, with $k$ the maximum lexicon
size (action words, priority senders, collaborators, skip lexicons).
The predicates are each $O(1)$ or an $O(k)$ substring/regex scan over a
bounded token blob; no cross-record work is done, so triage scales
linearly with the store and admits trivial parallelisation.

\subsubsection{The signal set}

\begin{table*}[t]
\centering
\small
\caption{The ten urgency signal predicates. Default weights sum to
$11$ (non-age) plus at most $3$ (age tier), giving the
$U(r) \in [0, 14]$ contract at defaults. Weights are per-key
overridable; unknown keys warn, non-integers raise.}
\label{tab:signals}
\begin{tabular}{lcp{9.2cm}}
\toprule
\textbf{Signal} & \textbf{Default $w$} & \textbf{Predicate / rationale} \\
\midrule
\texttt{mention}         & 3 & Regex \verb|@<you>\b| matches in \texttt{r.text} (case-insensitive). Empty \texttt{cfg.you} is guarded to avoid collapsing the regex to \verb|@\b|, which would fire on every raw \texttt{@}. Direct addressing is the single strongest human signal of obligation. \\
\texttt{priority\_sender} & 2 & \texttt{r.from\_name.lower()} $\in$ \texttt{cfg.priority\_senders}. Configured per-user (managers, VIPs). Off by default (empty set); nothing about any specific person or organization is baked in. \\
\texttt{action\_word}    & 2 & Any token from the action lexicon (defaults: \texttt{please}, \texttt{can you}, \texttt{need}, \texttt{review}, \texttt{approve}, \texttt{decision}, \texttt{deadline}, \texttt{asap}, \texttt{urgent}, \texttt{waiting on}, \texttt{action required}, \dots) appears in \texttt{subject $\|$ text}. Language-level cue that a reply is expected. \\
\texttt{question}        & 1 & \texttt{'?'} in \texttt{r.text}. Weak cue but cheap and orthogonal to action words. \\
\texttt{unread}          & 1 & \texttt{r.unread} is true. Loader-supplied; degrades gracefully to $0$ for sources with no read state. \\
\texttt{no\_reply}       & 1 & \texttt{r.thread\_id} is set and \texttt{r.replied} is false --- the user has not yet responded in a live thread. \\
\texttt{collaborator}    & 1 & \texttt{r.from\_name.lower()} $\in$ \texttt{cfg.collaborators}. Nudges known peers up without over-boosting them. \\
\texttt{age\_7d}         & 3 & Record age $> 168$\,h. Mutually exclusive with the finer age tiers. Escalates neglected obligations. \\
\texttt{age\_48h}        & 2 & Record age $\in (48, 168]$\,h. \\
\texttt{age\_24h}        & 1 & Record age $\in (24, 48]$\,h. \\
\bottomrule
\end{tabular}
\end{table*}

Two properties of the signal set deserve emphasis. First, the signals
are \emph{orthogonal by design}: \texttt{mention} is content-derived,
\texttt{priority\_sender} is identity-derived, \texttt{unread} is
platform-state-derived, and the age tiers are time-derived. A single
record is expected to fire two to four signals; the arithmetic sum is
then a joint score across independent evidence axes. Second, all
identity-bearing lexicons --- priority senders, collaborators, action
words, skip senders/subjects --- default either to an empty set or to a
generic vocabulary (mail-bot names, ``weekly digest'' subjects). No
individual person or organization is hard-coded. This is what makes the
scorer portable across users without a preprocessing step and safe to
ship as an open-source default.

\subsubsection{Design contrast: rule-based vs.\ learned}

The prevailing academic and industrial approach to inbox
prioritisation is a supervised classifier ---
\cite{TODO-horvitz-priorities} in the classical setting, and more
recently BERT-style text encoders or LLM zero-shot rankers fronted by a
managed API. Both classes buy accuracy at three costs \coviber{}
declines to pay:

\begin{enumerate}
\item \textbf{Labeled data.} A classifier needs per-user labels
  (\emph{this was urgent}, \emph{this was noise}) to be well-calibrated
  for a specific mailbox. Rule-based scoring calibrates via
  configuration, which the user writes once and can inspect.
\item \textbf{Cloud egress.} Zero-shot LLM ranking exports the entire
  inbox contents to a hosted model. \coviber{}'s target deployment is
  local-first over sensitive corpora (\S\ref{sec:threat-model}); this
  is a non-starter.
\item \textbf{Opacity.} A neural score is a scalar with no
  human-legible provenance. Definition~\ref{def:urgency} returns
  $(U(r), S(r))$: every triage decision ships with the list of signals
  that fired and their individual weights, which is what makes the FYI
  regression above visible in the first place.
\end{enumerate}

The tradeoff is real: rule-based scoring cannot learn user-specific
idioms without configuration edits, and it will underrank a
subtly-urgent record that fires zero signals. The system is honest
about this: the persona layer (\S\ref{sec:persona}) and the MCP recall
surface (\S\ref{sec:mcp}) are how a downstream LLM client re-ranks or
overrides triage output on demand. Urgency is the fast, transparent
first pass; the LLM is the slow, contextual second pass.


% Verified: False
% Refutation notes:
%   Two factual errors misrepresent what the shipped code does:
%   
%   (1) TUPLE ARITY IS WRONG. The draft states `V(M)` is an "8-tuple" but the equation itself lists **nine** components: $(n, \overline{w}, \phi, \varepsilon, \rho_?, \rho_!, O, C, W)$. The code (`coviber/persona.py`, VoiceProfile dataclass, lines 28–38) has nine fields matching that list one-for-one (`n_messages`, `avg_words`, `formality`, `emoji_rate`, `question_rate`, `exclaim_rate`, `top_openers`, `top_closers`, `vocabulary`). It is a 9-tuple, not an 8-tuple. This is an internal inconsistency and misstates the code.
%   
%   (2) EMOJI REGEX RANGE IS WRONG. The draft claims the second character-class range covers "the Miscellaneous-Technical through Geometric-Shapes range (U+2300–U+25FF)". The code (`persona.py` line 25) is `re.compile("[\\U0001F300-\\U0001FAFF⌀-➿]")`. The literals `⌀` and `➿` are U+2300 and U+27BF respectively — the actual upper bound is **U+27BF**, not U+25FF. U+25FF is the end of Geometric Shapes; U+27BF is inside the Dingbats block, well beyond Geometric Shapes and past Miscellaneous Symbols (U+2600–U+26FF) and most of Dingbats (U+2700–U+27BF). Both the numeric endpoint and the block-name description ("Geometric-Shapes") misstate what the code actually matches. The comment at lines 21–24 in the source also references dingbats explicitly, confirming the intent is the wider Dingbats-inclusive range.
%   
%   Other technical claims verified against `coviber/persona.py`: formality formula $\phi = \text{clip}(0.5 + \Delta/4n)$ matches line 102; formal/casual marker vocabularies match lines 71–72 exactly; thresholds $\phi > 0.6$ formal, $\phi < 0.35$ casual, $\varepsilon > 0.15$ emoji match lines 53, 56; opener-vs-closer single-line disambiguation (`len(stripped_lines) > 1`) matches line 80; word regex `[A-Za-z][A-Za-z'-]+` matches line 19; vocab filters (stopword, len > 3, no apostrophe) match line 94; top-3/top-20 truncations match lines 110–112; system-prompt top-12 signature words matches line 60; empty-corpus sentinel wording matches lines 49–52; inference-free / local-only claim is upheld by the source (regex + Counter only, no imports of any model/embedding library).
% Notes to assembler: Section uses \path{...} (package: url), \operatorname{clip}, and \lceil/\rceil — all covered by amsmath + url which the preamble should already include. Uses \coviber{} macro (define as \newcommand{\coviber}{\textsc{CoViber}} in preamble, or global-replace before compile).

Minor deviation from prom

\subsection{Statistical persona modelling}
\label{sec:persona}

A common failure mode of LLM-assisted drafting is that the model,
lacking any knowledge of \emph{how a specific user writes}, defaults to a
generic assistant register. The obvious remedy --- fine-tune, or feed the
LLM a large corpus of the user's prior writing at each request --- is
either expensive or leaks the user's private history into a third-party
inference call. \coviber{} instead computes a compact statistical
\emph{voice profile} from the user's self-authored records and emits a
short natural-language \emph{system prompt} that any downstream LLM can
consume. \textbf{No LLM inference is performed anywhere in this
pipeline}; every quantity in the profile is derived by regular
expressions and counters over local text (\path{coviber/persona.py}).

\paragraph{The voice profile.}
Let $M = (m_1, \ldots, m_n)$ be the sequence of non-empty text bodies of
records whose \texttt{from\_name} matches the configured user identity.
The \texttt{VoiceProfile} learned from $M$ is the 8-tuple
\begin{equation}
  V(M) \;=\;
  \bigl(
    n,\;
    \overline{w},\;
    \phi,\;
    \varepsilon,\;
    \rho_?,\;
    \rho_!,\;
    O,\;
    C,\;
    W
  \bigr)
  \label{eq:voice}
\end{equation}
where $n$ is the number of messages, $\overline{w}$ is the mean word
count per message, $\phi \in [0,1]$ is the formality index defined
below, $\varepsilon$ is the mean emoji count per message,
$\rho_?, \rho_!$ are the mean per-message counts of \texttt{?} and
\texttt{!}, and $O, C, W$ are the multisets of the top-3 openers, top-3
closers and top-20 signature vocabulary items (each stored with a count).

\paragraph{Formality signal.}
Let \textsc{formal} and \textsc{casual} be fixed marker vocabularies:
\begin{align*}
  \textsc{formal} &= \{\text{\emph{please, kindly, regards, sincerely, would you, could you, i would}}\} \\
  \textsc{casual} &= \{\text{\emph{hey, yeah, gonna, lol, thx, cool, sure thing, np}}\}
\end{align*}
For each message $m_i$, let $f_i$ (resp.\ $c_i$) be the number of
substring occurrences of \textsc{formal} (resp.\ \textsc{casual})
markers in the lowercased body. Let
$\Delta = \sum_{i=1}^{n} (f_i - c_i)$. The formality index is
\begin{equation}
  \phi \;=\; \operatorname{clip}_{[0,1]}\!\left(\tfrac{1}{2} + \tfrac{\Delta}{4n}\right)
  \label{eq:formality}
\end{equation}
so that $\phi = 0.5$ is the neutral midpoint reached when
$\Delta = 0$, and the corpus must average one net marker per message
across four messages to move $\phi$ by $0.25$. The scale factor $4n$
makes $\phi$ robust to short bursty samples --- a single very formal
message cannot alone push a corpus to $\phi \to 1$.

\paragraph{Opener/closer detection.}
Naively regex-matching openings and closings on the whole message body
produces double counts on the very common \emph{single-line}
acknowledgement (``\texttt{Thanks!}''), because it matches both the
opener regex (\texttt{thanks} $\in$ greeting alternation) and the closer
regex. \coviber{} avoids this by treating single-line messages as
having an opener \emph{or} a closer, not both. Concretely, let
$L(m_i) = (\ell_1, \ldots, \ell_k)$ be the non-empty lines of $m_i$.
Then the opener regex is matched against $\ell_1$; the closer regex is
matched against $\ell_k$ only when $k > 1$
(\path{persona.py}, audit L6/\#26). This makes $|O|$ and $|C|$
independently meaningful signals.

\paragraph{Signature vocabulary filter.}
The vocabulary counter $W$ ranks tokens produced by the regex
\texttt{[A-Za-z][A-Za-z'-]+} against three filters, applied in order:
(i) reject a fixed stopword list of pronouns, articles, and function
words; (ii) reject tokens of length $\leq 3$; (iii) reject any token
containing an ASCII apostrophe. Filter (iii) is deliberate: contractions
such as \emph{don't}, \emph{I'll}, \emph{we're} are essentially uniform
across English speakers and carry no idiolectal signal, so admitting
them would dominate the top-$k$ list and drown out the actual signature
terms (audit L6/\#28). The result is a compact ordered list of at most
20 distinctive multi-character words.

\paragraph{Emoji regex.}
The emoji counter uses a single compiled character-class regex covering
the Unicode Supplementary Multilingual Plane emoji blocks
(\texttt{U+1F300}--\texttt{U+1FAFF}) plus the Miscellaneous-Technical
through Geometric-Shapes range (\texttt{U+2300}--\texttt{U+25FF}); the
latter picks up clock, watch, and simple-shape glyphs that occur in
professional writing and would otherwise be missed (audit L6/\#30).

\paragraph{System prompt.}
For a display name $\eta$, \texttt{VoiceProfile.system\_prompt}$(\eta)$
returns a single-paragraph natural-language instruction of the form
\begin{quote}\small\itshape
Write as $\eta$. Voice: $\{$formal $\mid$ neutral $\mid$ casual$\}$,
$\sim\!\!\lceil\overline{w}\rceil$ words per message,
$\{$occasionally $\mid$ rarely$\}$ uses emoji. Typically opens with
``$o_1$'' and closes with ``$c_1$''. Signature vocabulary:
$w_1, \ldots, w_{12}$. Match this tone and length; do not over-explain.
\end{quote}
with the register bucketed by thresholds $\phi > 0.6$ (formal),
$\phi < 0.35$ (casual), else neutral; the emoji clause switches at
$\varepsilon > 0.15$ emojis per message. This prompt is the sole
artefact consumed by the downstream LLM --- there is no embedding
lookup, no fine-tune, no per-message inference on \coviber{}'s side.

\paragraph{Empty-corpus sentinel.}
When \texttt{n\_messages} $= 0$, the system prompt does not fabricate
plausible-sounding defaults (which would encode a spurious ``formal, 0
words per message'' voice) but returns the sentinel string
\begin{quote}\small\itshape
No self-authored writing samples were found for $\eta$. No voice
profile learned --- draft in your own default voice.
\end{quote}
This is load-bearing: an MCP client receiving this string can transparently
fall back to its own defaults without producing a corrupted persona
(audit L6/\#29).

\paragraph{Complexity and privacy.}
Learning $V(M)$ is a single linear pass over the concatenated text with
constant per-token work, so $|V(M)| = O(\sum_i |m_i|)$ time and
$O(|W| + |O| + |C|)$ space --- bounded by the top-$k$ truncations, i.e.\
$O(1)$ in the corpus size. Because every operation is a local regex or
counter increment, the profile can be recomputed on every
\texttt{voice\_profile} MCP call (\S\ref{sec:mcp}) without cache
invalidation logic; the only I/O is the local record store read.
Crucially, no message body ever leaves the host during profile
construction --- the profile is derived, emitted, and discarded on the
same machine.

% Verified: False
% Refutation notes:
%   Two claims are contradicted by the shipped code:
%   
%   (1) HARD REFUTATION — Protocol method count. The draft states in §3.7.2 that the VectorStore protocol is "intentionally small — six methods". The actual Protocol in /Users/I869061/code/active/coviber-oss/coviber/vector_stores.py (lines 25-45) defines SEVEN methods: known_ids, upsert, delete, search, model, rebind, wipe. The section's own Listing~\ref{lst:vs-proto} also enumerates all seven. The prose "six methods" contradicts both the code and the listing shown in the same section. Fix: change "six methods" to "seven methods".
%   
%   (2) Softer inaccuracy — model-change invalidation for JSONVectorStore. The draft claims "Every backend persists its embedder tag ... so that a model change is detected on the first call after startup and triggers wipe(), forcing a full re-encode." This is only true for QdrantVectorStore (vector_stores.py lines 224-232 explicitly call self.wipe() in _reconcile_model_tag on mismatch). JSONVectorStore does NOT call wipe() on model mismatch — its _load() (lines 76-78) simply returns an empty dict on mismatch, and the next _persist() overwrites embeddings.json with the new tag. The full re-encode does happen, but via lazy cache invalidation rather than a wipe() call. Fix: describe the invalidation mechanism more generically ("triggers a lazy invalidation / rebuild" or "invalidates the cached vectors") or split the sentence between the two backends.
%   
%   Minor notational issue (not sufficient alone to refute, but flagging): the keyword score's denominator is documented as log(2 + |r.blob|) where |r.blob| is ambiguous; the code uses log(2 + len(blob.split())), i.e. word count over the concatenated subject+text+from_name string. If |r.blob| is intended as "length in words" this is fine; a footnote making that explicit would remove the ambiguity.
% Notes to assembler: Section uses two forward references: \\ref{sec:design-principles} (for the local-first stance and the zero-dependency-by-default principle) and \\S\\ref{sec:record-store} etc. within-section labels. Confirm the design-principles section is labeled sec:design-principles or adjust.\n\nPreamble require

\subsection{Persistence and semantic search}
\label{sec:persistence-search}

CoViber's persistence layer is designed around three assumptions that
follow from its local-first stance (\S\ref{sec:design-principles}):
records must survive process crashes with no torn state; concurrent
writers (a long-running MCP server and an on-demand CLI \texttt{ingest})
must serialize without a database daemon; and the semantic-search index
must be pluggable so a user can start on a zero-dependency JSON file and
migrate to Qdrant once the corpus grows past what fits comfortably in
memory. This section describes the record store
(\S\ref{sec:record-store}), the pluggable vector-store abstraction and
its two shipped backends (\S\ref{sec:vector-stores}), the delta-encoding
invariant that keeps the warm search path $O(\Delta)$
(\S\ref{sec:delta-invariant}), and the keyword-scoring fallback that
preserves usable behaviour when the optional embedding extra is absent
(\S\ref{sec:keyword-fallback}).

\subsubsection{The record store}
\label{sec:record-store}

Records are persisted as one JSON object per line in
\texttt{<data\_dir>/records.jsonl}. Each ingestion cycle is an
atomic\hyp{}rewrite: \texttt{Store.upsert} reads the full file, merges
the new records into a dict keyed by \texttt{Record.id}
(the content-derived MD5 hash minted in
\texttt{Record.\_\_post\_init\_\_}, so two loaders that produce the same
$(\texttt{text},\texttt{subject},\texttt{from\_name},\texttt{source})$
tuple deduplicate naturally), rewrites the merged set to a temporary
file, and finally invokes \texttt{os.replace} to swap it into place.
The temporary filename is tagged with the writer's PID and a random
suffix (\texttt{records.jsonl.tmp.\{pid\}.\{uuid\}}), so two writers
that somehow bypass the advisory lock — a future bulk-import helper, a
platform with neither \texttt{fcntl} nor \texttt{msvcrt} — cannot
collide on a shared temporary path.

\paragraph{Advisory lock.}
The read–merge–write cycle is enclosed in an advisory file lock on
\texttt{<data\_dir>/.lock}, acquired via \texttt{fcntl.flock} on POSIX
and \texttt{msvcrt.locking} on Windows. The lock deliberately covers
the read as well as the write: were the lock scoped only around the
final rename, two concurrent writers could each read the pre-merge
file, each merge their disjoint updates, and the last writer's rename
would silently overwrite the first writer's contribution. Readers do
not take the lock at all — the atomic \texttt{os.replace} in
\texttt{upsert} already guarantees that a reader observes either the
pre-write or post-write file, never a torn one. This asymmetry —
writers serialize, readers proceed lock-free — is deliberate: the MCP
tools (\texttt{recall}, \texttt{catch\_me\_up}, \texttt{who\_is}) run
on the hot path of an LLM turn and must not block on unrelated
ingestion. On platforms exposing neither \texttt{fcntl} nor
\texttt{msvcrt}, the context manager degrades to a no-op; single-process
correctness is preserved, only the multi-writer guarantee is lost, and
this is documented at the call site.

\paragraph{Corrupt-line quarantine.}
For a memory product, silent data loss is the worst failure mode. When
\texttt{\_read\_all} encounters a line that fails either
\texttt{json.loads} or \texttt{Record.from\_dict}, it does not raise
(which would brick every subsequent \texttt{upsert}) and it does not
drop the line (which would erase evidence on the next rewrite).
Instead the raw line is appended verbatim to
\texttt{records.jsonl.bad} and a warning is emitted to \texttt{stderr}
with the file and line number. On the next \texttt{upsert} the
quarantined line is thus present in the sidecar but absent from
\texttt{records.jsonl}, so the corpus continues to grow while the
operator has a durable record of every parsing failure.

\paragraph{Forward-compatible schema drops.}
\texttt{Record.from\_dict} tolerates unknown keys — an older coviber
reading a store written by a newer version does not crash. But the
first time a given unknown key is observed in a process,
\texttt{Record.from\_dict} emits a \texttt{RuntimeWarning} identifying
the dropped field, so an operator running an accidental downgrade sees
that the next \texttt{upsert} rewrite will lose those fields. Per-key
one-shot deduplication prevents an $N$-record store from producing $N$
identical warnings.

\subsubsection{Pluggable vector stores}
\label{sec:vector-stores}

Semantic search is orthogonal to record persistence: the same
records.jsonl can be indexed by any backend that satisfies the
\texttt{VectorStore} protocol (Listing~\ref{lst:vs-proto}). The
protocol is intentionally small — six methods — because it is the
seam between coviber's zero-dependency default path and every future
vector-database integration.

\begin{lstlisting}[float,caption={The VectorStore protocol
(\texttt{coviber/vector\_stores.py}).},label={lst:vs-proto},language=Python,basicstyle=\small\ttfamily]
class VectorStore(Protocol):
    def known_ids(self) -> set[str]: ...
    def upsert(self, id_vecs: list[tuple[str, list[float]]]) -> None: ...
    def delete(self, ids: Iterable[str]) -> None: ...
    def search(self, query_vec: list[float],
               live_ids: set[str], limit: int
               ) -> list[tuple[float, str]]: ...
    def model(self) -> str: ...
    def rebind(self, model: str) -> None: ...
    def wipe(self) -> None: ...
\end{lstlisting}

The \texttt{model()} accessor is load-bearing: two vectors produced by
different sentence encoders inhabit different metric spaces and cannot
be meaningfully mixed. Every backend persists its embedder tag
(\texttt{all-MiniLM-L6-v2} in the shipped default) so that a model
change is detected on the first call after startup and triggers
\texttt{wipe()}, forcing a full re-encode. \texttt{search} takes the
set of currently-live record ids and filters against it inside the
backend, so a record removed from \texttt{records.jsonl} but not yet
purged from the index cannot leak into results during a partial
reconciliation.

\paragraph{JSONVectorStore --- the default.}
The zero-dependency backend stores vectors in a single file
\texttt{embeddings.json} alongside \texttt{records.jsonl}, with
schema $\{\texttt{model}: \tau,\, \texttt{vectors}: \{r_i \mapsto \mathbf{v}_i\}\}$.
It has no in-process cache: every call re-reads the file, matching
pre-v0.3 semantics and letting an external process edit the index
without a coviber restart. Vectors are coerced to
\texttt{float32}-precision Python floats before serialization; this
roughly halves on-disk width relative to \texttt{float64} with no
meaningful precision loss for a \texttt{float32}-native encoder. Writes
use the same PID\hyp{}and\hyp{}uuid tagged temporary path idiom as the
record store, so a lock-free \texttt{search()} on the read path cannot
race a \texttt{upsert()} into a torn file. A cache whose model tag
does not match the current tag is discarded silently and rebuilt
lazily by the next search — the same treatment as an unreadable or
non-dict cache file. The JSON backend is documented to be appropriate
up to $\sim 10^{5}$ records; past that, the whole-file read on every
call becomes the dominant cost.

\paragraph{QdrantVectorStore --- the ANN escape hatch.}
Beyond that scale the shipped Qdrant\cite{TODO-qdrant} backend
delegates persistence and search to a server-side ANN index. It is
opt-in behind the \texttt{[qdrant]} extra, so the default install
still ships with a hard dependency count of zero (see
\S\ref{sec:design-principles}). Two implementation choices are worth
noting. First, coviber's \texttt{Record.id} is exactly a 128-bit MD5
digest, which parses into a UUID with no loss of bits and gives a
stable, invertible mapping from record id to Qdrant point id — no
external id-to-uuid table is required. Second, the model tag is
stored in a small sidecar \texttt{qdrant.meta.json} on the local
filesystem rather than in Qdrant collection metadata, so
model-invalidation logic is identical to the JSON backend and does not
depend on Qdrant client capabilities that may differ across versions.

\paragraph{Resolver and the local-first invariant.}
The \texttt{resolve()} function is the single place in the codebase
that maps configuration to a backend instance. Its precedence is
(1)~an explicit \texttt{qdrant:} block in \texttt{coviber.yaml}, then
(2)~the \texttt{COVIBER\_QDRANT\_URL} environment variable, then
(3)~\texttt{JSONVectorStore}. When a Qdrant URL is configured but its
host does not textually match loopback (\texttt{localhost},
\texttt{127.x.x.x}, \texttt{::1}), the resolver emits a
\texttt{stderr} warning that vectors will be transmitted off-device.
This is intentionally non-fatal — a remote Qdrant is a legitimate
power-user configuration — but must be visible: a copy-pasted managed
cluster URL should not silently violate the local-first invariant.
Correspondingly, a Qdrant URL that fails on connect raises up to the
caller rather than falling back to JSON, because a silent fallback
would mask real configuration drift for a system whose sole purpose is
durable memory.

\subsubsection{The delta-encoding invariant}
\label{sec:delta-invariant}

Semantic search proceeds against the union of the record store and the
vector store, but re-encoding the whole corpus on every query would
make the warm path scale with corpus size $N$ rather than with the
number of newly ingested records. Algorithm~\ref{alg:delta} states the
invariant coviber's \texttt{Store.\_embedding\_search} enforces on
every call.

\begin{algorithm}[t]
\caption{Delta-reconcile-and-search (\texttt{Store.\_embedding\_search}).}
\label{alg:delta}
\begin{algorithmic}[1]
\Require query string $q$, record set $\mathcal{R}$, vector store $\mathcal{V}$, encoder $E$, top-$k$
\State $L \gets \{r.\mathrm{id} : r \in \mathcal{R}\}$ \Comment{live ids}
\State $K \gets \mathcal{V}.\Call{known\_ids}{}$
\State $S \gets K \setminus L$ \Comment{stale: indexed but no longer in store}
\State $M \gets \{r \in \mathcal{R} : r.\mathrm{id} \notin K\}$ \Comment{missing: in store, unindexed}
\If{$M \neq \emptyset$}
    \State $\mathbf{V} \gets E.\Call{encode}{\{r.\mathrm{subject} \parallel r.\mathrm{text} : r \in M\}}$
    \State $\mathcal{V}.\Call{upsert}{\{(r.\mathrm{id}, \mathrm{f32}(\mathbf{v}_r)) : r \in M\}}$
\EndIf
\If{$S \neq \emptyset$} \State $\mathcal{V}.\Call{delete}{S}$ \EndIf
\State $\mathbf{q} \gets E.\Call{encode}{q}$
\State \Return $\mathcal{V}.\Call{search}{\mathbf{q},\ L,\ \max(2k, 16)}$ truncated to $k$
\end{algorithmic}
\end{algorithm}

\paragraph{Complexity.}
Let $N = |\mathcal{R}|$ and $\Delta = |M| + |S|$ be the number of
records changed since the last search. The set operations
$K \setminus L$ and the membership filter for $M$ are both $O(N)$ in
Python's hashed set implementation, but the encoder call — the
dominant cost — runs on exactly $|M|$ inputs. On a warm store where
ingestion has added $|M|$ records and removed $|S|$, the reconciliation
step is $O(\Delta)$ in the expensive dimension (embedder invocation and
on-disk vector I/O). On the cold path where $K = \emptyset$ we
recover the trivial $O(N)$ full re-encode. The search step itself is
$O(N)$ for \texttt{JSONVectorStore} (dot product against every live
vector) and asymptotically sublinear for
\texttt{QdrantVectorStore} (server-side HNSW\cite{TODO-hnsw}); this
motivates the $10^{5}$-record migration threshold.

\paragraph{Over-fetch and live-id filtering.}
The final \texttt{search} call over-fetches to $\max(2k, 16)$ hits and
truncates to $k$ after mapping ids back to \texttt{Record} objects.
This absorbs two edge cases without additional lock discipline:
a record that was deleted from \texttt{records.jsonl} between the
\texttt{known\_ids()} call and the search cannot inflate the result
above $k$ live entries, and a hit whose id no longer maps to a live
\texttt{Record} is dropped by the reverse-lookup filter in
\texttt{\_embedding\_search}.

\paragraph{Why not re-encode on ingest?}
An alternative design would encode inside \texttt{upsert} and skip
reconciliation on the search path. Coviber deliberately does not do
this: the embedder is an optional dependency, so \texttt{upsert} must
work when \texttt{sentence-transformers} is not installed;
\texttt{ingest} is called from environments (CI, cron) where a
1--2\,GB model download is undesirable; and encoding on the search
path lets a user install the \texttt{[search]} extra later without
re-ingesting. The delta invariant makes deferred encoding cheap
enough that the design cost is paid only on the first search after
each ingest.

\subsubsection{The keyword fallback}
\label{sec:keyword-fallback}

When \texttt{sentence-transformers} is not installed —
\texttt{\_get\_embedder} raises \texttt{ImportError} —
\texttt{Store.search} catches it \emph{silently} and falls back to the
keyword-scoring path:
\[
\mathrm{score}(r, q) = \frac{\sum_{t \in q} \mathrm{count}(t,\, r.\text{subject} \parallel r.\text{text} \parallel r.\text{from\_name})}{\log\bigl(2 + |r.\text{blob}|\bigr)},
\]
records with zero matching terms are pruned, and the remainder are
returned sorted by descending score. This is a term-frequency scoring
scheme — deliberately not TF--IDF, since an IDF pass would either
require another persisted index or an $O(N)$ scan on every query
where the embedding path pays only $O(\Delta)$. The logarithmic
length normalization prevents very long records from dominating on
raw term counts alone.

Two operationally meaningful details distinguish this fallback from a
naive downgrade. First, the \texttt{ImportError} branch is silent
because it represents an expected configuration; the
\texttt{last\_search\_backend} attribute records
\texttt{"keyword"} so callers (and MCP tools) can report the actual
backend to the user. Second, an unexpected exception during embedding
search — CUDA OOM, HuggingFace-hub transient failure on first model
download, a torn embedding cache, an unreachable Qdrant — is
\emph{surfaced} to \texttt{stderr} and via a \texttt{RuntimeWarning},
even though the same keyword fallback is invoked. This lets an
operator distinguish ``extras missing'' (silent, expected) from
``extras installed but broken'' (loud, diagnostic), a distinction that
matters when a search result set suddenly looks worse without a
corresponding configuration change.


% Verified: True
% Notes to assembler: Section uses \subsection{...} and \label{sec:mcp} — assembler should verify parent \section{} (System Design, §3) exists and adjust nesting if needed.

Cross-references used: \ref{sec:ingestion}, \ref{sec:workgraph}, \ref{sec:urgency}, \ref{sec:persona}. Ensure these labels match what the sibling se

\subsection{The MCP-Native Interface}
\label{sec:mcp}

The Model Context Protocol (MCP) is a JSON-RPC framing specification for exposing tools, resources, and prompts to a host LLM client~\cite{TODO-mcp-spec}. A conformant \emph{server} advertises a schema-typed set of tools; a conformant \emph{client} (Claude Desktop, Claude Code, or any third-party host that implements the same specification) discovers those tools at startup and marshals structured calls into and out of the model. The protocol deliberately separates the memory substrate from the reasoning agent: a memory engine that speaks MCP inherits every current and future client for free, and, symmetrically, an agent that speaks MCP inherits every memory engine.

This separation is the design commitment of CoViber\Hyphdash{}\textsc{oss}. The pipeline of \S\ref{sec:ingestion}\Hyphdash\S\ref{sec:urgency} produces a purely local artifact set --- a deduped record log, a work graph, and (optionally) a vector index --- with no direct client coupling. The MCP server is a thin adapter that projects those artifacts into the seven tool signatures listed in Table~\ref{tab:mcp-tools}. Every capability the whitepaper describes is reachable through this surface; there is no privileged client and no ``native'' UI.

\begin{table}[t]
\small
\centering
\begin{tabular}{@{}p{0.28\linewidth}p{0.66\linewidth}@{}}
\toprule
\textbf{Tool} & \textbf{Semantics} \\
\midrule
\texttt{recall(query, limit=6)} & Semantic (or keyword-fallback) search over the local record log, returning the top-$k$ records by cosine similarity to the query. \\
\texttt{catch\_me\_up(limit=10)} & The urgency-ranked action queue $Q$ of \S\ref{sec:urgency}, truncated to \texttt{limit} highest-scoring records with their firing signals. \\
\texttt{who\_is(name)} & The person node from the work graph (\S\ref{sec:workgraph}); case-insensitive on \texttt{name}, returns platforms, projects, channels, and interaction count. \\
\texttt{project\_status(name)} & The project node from the work graph: participants, channels, and cross-referenced tickets. \\
\texttt{graph\_summary()} & Aggregate counts (\texttt{people}, \texttt{projects}, \texttt{channels}, \texttt{tickets}), the enumerated project list, and the top eight most-interactive people. \\
\texttt{voice\_profile()} & The statistical \texttt{VoiceProfile} of \S\ref{sec:persona} learned over self-authored records, paired with a ready-to-paste \texttt{system\_prompt(you)}. \\
\texttt{refresh(loader, path="")} & Trigger one \(\text{load} \to \text{store} \to \text{graph}\) cycle through any registered loader; returns a stats string. \\
\bottomrule
\end{tabular}
\caption{The MCP tool surface exposed by \texttt{coviber.mcp\_server}. Each tool is a single-purpose projection of a pipeline artifact; no tool composes state across calls.}
\label{tab:mcp-tools}
\end{table}

\paragraph{Configuration parity and hot reload.}
The MCP server and the command-line entry point share one settings parser (\texttt{coviber.config.read\_config}), so an operator can point \texttt{coviber triage --config foo.yaml} and \texttt{coviber serve --config foo.yaml} at the same file with identical semantics. Crucially, the server \emph{re-reads} the config on every tool invocation via a private \texttt{\_settings()} helper: an edit to the YAML file --- adding a priority sender, adjusting an urgency weight, switching the loader --- takes effect on the next call without restarting the process or reattaching the client. Environment variables \texttt{COVIBER\_CONFIG}, \texttt{COVIBER\_DATA\_DIR}, and \texttt{COVIBER\_YOU} act as overrides, evaluated in the same order at every call.

\paragraph{Defensive graph reads.}
Three tools (\texttt{who\_is}, \texttt{project\_status}, \texttt{graph\_summary}) consume \texttt{workgraph.json}. The read path is centralized in \texttt{Store.load\_graph}, which returns a \texttt{dict} on success and a human-readable string sentinel on either of two failure modes: (a) the file does not exist yet --- i.e.\ the operator has not run \texttt{coviber ingest} --- or (b) the file exists but fails to parse (a legacy layout, a hand-edit, external corruption). Each MCP tool inspects the return type and forwards the sentinel verbatim, so a torn graph produces a diagnostic message on the client rather than a Python traceback across the JSON-RPC boundary. All internal accesses use \texttt{dict.get} with defaults, so an incomplete graph (missing \texttt{tickets}, older schema) still yields a valid summary.

\paragraph{Loader errors surface as strings.}
The \texttt{refresh} tool catches the enumerated set \{\texttt{KeyError}, \texttt{ValueError}, \texttt{FileNotFoundError}, \texttt{IsADirectoryError}, \texttt{PermissionError}, \texttt{OSError}\} and converts each to a one-line error string of the form \texttt{"refresh failed ($T$: $\dots$)"}. \texttt{KeyError} covers an unknown loader name; the remaining exceptions cover the file-system paths a loader can plausibly touch. \texttt{ImportError} is deliberately \emph{not} caught: it indicates the operator requested a loader gated on an optional extra (\texttt{[scrape]}, \texttt{[qdrant]}) without installing it, which is a setup problem worth surfacing plainly. All other exceptions propagate to FastMCP, which converts them to a per-tool error without killing the server.

\paragraph{Transport and privacy invariant.}
The server binds only to standard input and output; there is no listening socket, and one server process serves at most one client at a time. Combined with the local-first pipeline of \S\ref{sec:ingestion}, this yields a structural privacy invariant: the only network traffic issued by CoViber\Hyphdash{}\textsc{oss} is whatever the \texttt{[scrape]} extra (an operator-triggered web fetch inside a \texttt{refresh} call) or the \texttt{[qdrant]} extra (a connection to a Qdrant instance the operator configured) emits at the operator's explicit request. There is no telemetry channel, no update check, and no upstream inference call --- the persona module is inference-free by construction (\S\ref{sec:persona}), and the embedding model, when enabled, runs locally through \texttt{sentence-transformers}.

\paragraph{Deployment.}
A client's configuration --- e.g.\ Claude Desktop's \texttt{claude\_desktop\_config.json} --- registers the server as a stdio command, \texttt{coviber serve}, optionally with \texttt{--config} and \texttt{--data-dir} flags. The client spawns one server per session, discovers the seven tools, and thereafter treats them like any other MCP capability. The same server binary is consumed by Claude Code from the command line and, in principle, by any third-party MCP host: the memory engine is model- and vendor-agnostic by construction.


% -----------------------------------------------------------------------
%  §4  Evaluation
% -----------------------------------------------------------------------
% Verified: False
% Refutation notes:
%   §4.5 (Persona-profile determinism) makes an empirical claim about the demo corpus that is contradicted by the shipped code. Specifically:
%   
%   1. "Empirically, on the demo corpus's three authored messages the top opener ('Hi'), top closer ('Thanks'), and top twelve vocabulary items each have strict count leads, so the system_prompt output ... is identical across all 3! = 6 permutations of the input."
%   
%   Tracing persona.learn against demo.py's three from_name=="you" messages:
%   
%   - Msg 1 first line "Hi Grace, reviewed the plan — looks good to ship. One ask: let's gate the rollout on the canary metrics. Thanks, " → _GREETING matches "hi" → openers["Hi"] += 1.
%   - Msg 2 first line "Thanks for flagging — looking now. Can you drop the trace id? Will patch today." → _GREETING matches "thanks" (it is one of the greeting alternatives in `r"^\s*(hi|hey|hello|dear|good morning|good afternoon|thanks|thank you|team)\b"`) → openers["Thanks"] += 1.
%   - Msg 3 first line "Nice approach on the backoff. ..." → no match.
%   
%   Result: openers = Counter({"Hi": 1, "Thanks": 1}) — a tie, NOT a strict count lead. Under Counter.most_common's insertion-order tie-break, top_openers[0][0] is "Hi" for orderings that see msg 1 before msg 2, and "Thanks" for the reverse orderings. The system_prompt reads exactly top_openers[0][0], so it is NOT identical across all 6 permutations — it changes depending on which authored message is encountered first.
%   
%   2. The closer claim is even more directly wrong. Per persona.py's L6/#26 fix, `last = stripped_lines[-1] if len(stripped_lines) > 1 else ""` — a single-line message contributes no closer. All three authored demo messages are single-line (no `\n`), so closers = Counter() (empty). top_closers is empty, and system_prompt uses its hardcoded fallback `cl = "Thanks"` (persona.py line 55: `cl = self.top_closers[0][0] if self.top_closers else "Thanks"`). The draft attributes the "Thanks" in the output to "top closer" and claims a strict count lead there — neither is true. The value is constant only because it's a fallback constant, not because a strict lead exists in top_closers.
%   
%   3. The draft's own assembler note correctly claims "vocabulary is 3/7/4 and 1.00 by construction" for the work-graph table (§4.4), which I verified against the code and demo corpus — that part is fine. But the persona empirical paragraph in §4.5 makes a stronger statement than the code supports.
%   
%   The Proposition itself and the proof sketch are correct — that's a general statement about the algorithm. Only the specific empirical claim about the demo corpus is refuted. All other claims in the section (MD5 record ids, bench harness details including repeat counts of 3 vs 5, Store.last_search_backend behavior, JSON whole-file load per query, docstring quote about ~10^5 records, work-graph ground truth 3/7/4, project-match word-boundary, ticket normalization for "PR #482" variants, case-preserving people node display_name, cited audit findings L4/#11 and L4/#14, inference-free / counter-based persona) match the shipped code exactly.
%   
%   Recommended fix: either drop the "empirically ... strict count leads ... identical across all 3! = 6 permutations" sentence entirely, or reword it to reflect that (a) top_openers has a 1-1 tie between "Hi" and "Thanks" so the emitted opener is order-dependent under insertion-order tie-breaking, and (b) top_closers is empty on the demo corpus (single-line messages), so the "Thanks" that appears in system_prompt comes from the hardcoded fallback and is order-independent for that trivial reason.
% Notes to assembler: Preamble additions needed:
  - \\usepackage{siunitx}  (for \\SI{XXX}{\\milli\\second})
  - \\usepackage{booktabs} (for \\toprule/\\midrule/\\bottomrule)
  - \\usepackage{amsthm}   (for the \\begin{proposition} block — either define a proposition environment via \\newtheorem{proposition}{Proposition}

\section{Evaluation}
\label{sec:eval}

The evaluation has three goals: (i) establish that each stage of the pipeline
scales to the corpus sizes a single knowledge worker actually accumulates over
a career (order $10^{5}$ records), (ii) show that the pluggable vector backend
crosses over from the zero-dependency JSON store to a server-side ANN index at
a predictable, corpus-size boundary, and (iii) sanity-check the deterministic
components (work-graph extraction, persona statistics) on the ground-truth demo
corpus. All experiments run on the synthetic \texttt{demo} loader that ships with
the package (\S\ref{sec:synthetic-corpus}); no proprietary data is involved and
every number in this section is reproducible via a single command:
\begin{center}
\texttt{python bench/run.py -{}-scale <N>}
\end{center}

\subsection{Reproducibility harness}
\label{sec:bench-harness}

The harness in \texttt{bench/run.py} deterministically expands the fifteen-record
demo corpus to a target scale $N$ by cycling the templates and appending a
per-record index to the \texttt{subject} and \texttt{text} fields, so that record
identity (an MD5 over the salient fields, see \S\ref{sec:record-model}) remains
unique. It then times the four pipeline stages named in the whitepaper hot path,
each as the median of $k \in \{3,5\}$ repeats measured with
\texttt{time.perf\_counter}: ingest with deduplication into the JSONL store,
work-graph construction, full-scan urgency triage, and top-$8$ semantic search
against the pluggable vector backend. The search timing is deliberately the
\emph{warm} median: the first repeat pays the encoder-load and cache-population
cost, the subsequent repeats measure query-only latency, which is the metric
that matters for interactive MCP use. The backend that actually served the
query is recorded (\texttt{Store.last\_search\_backend}) so the reported number
is never ambiguous between the JSON path, the Qdrant path, and the keyword
fallback.

\subsection{End-to-end latency by stage}
\label{sec:eval-latency}

Table~\ref{tab:latency} reports median wall-clock latency for each hot-path
stage at three corpus scales. The numbers reflect the JSON vector backend
(\texttt{JSONVectorStore}) — the zero-dependency default — on a single
commodity laptop; the exact hardware fingerprint is written into the results
file described in \S\ref{sec:eval-repro}.

\begin{table}[h]
\centering
\small
\begin{tabular}{lrrr}
\toprule
\textbf{Stage} & \textbf{$N{=}2\,000$} & \textbf{$N{=}10^{4}$} & \textbf{$N{=}10^{5}$} \\
\midrule
Ingest + dedup (JSONL) & \SI{XXX}{\milli\second} & \SI{XXX}{\milli\second} & \SI{XXX}{\milli\second} \\
Work-graph build       & \SI{XXX}{\milli\second} & \SI{XXX}{\milli\second} & \SI{XXX}{\milli\second} \\
Urgency triage (full scan) & \SI{XXX}{\milli\second} & \SI{XXX}{\milli\second} & \SI{XXX}{\milli\second} \\
Semantic search (top-8, warm) & \SI{XXX}{\milli\second} & \SI{XXX}{\milli\second} & \SI{XXX}{\milli\second} \\
\bottomrule
\end{tabular}
\caption{Median latency per pipeline stage on the synthetic corpus, JSON
vector backend. Reproduce via \texttt{python bench/run.py -{}-scale N} on any
machine; the harness fingerprints the hardware for cross-run comparability.}
\label{tab:latency}
\end{table}

Three observations follow from the algorithmic structure of each stage, and
should hold in the measured numbers. Ingest performs a single read of
\texttt{records.jsonl}, an $O(|R_{\text{new}}|)$ dedup merge against the
in-memory id map, and one write-then-rename; wall time is I/O-bound and grows
approximately linearly with $N$. Work-graph construction is a single pass over
all records (\S\ref{sec:workgraph}) with $O(1)$ per-record work modulo regex
matching, again linear in $N$. Urgency triage is a full scan applying a fixed
number of signal predicates per record with $O(|A|)$ substring checks over the
configured action-word set $A$ (\S\ref{sec:urgency}); linear in $N \cdot |A|$.
Warm semantic search under \texttt{JSONVectorStore} loads the entire
persisted vector map on every call (\S\ref{sec:vector-stores}); the query is
$O(N \cdot d)$ for embedding dimension $d = 384$. This last fact motivates
the Qdrant backend and the crossover experiment in
\S\ref{sec:eval-jsonvsqdrant}.

\subsection{JSON versus Qdrant crossover}
\label{sec:eval-jsonvsqdrant}

The two vector backends serve identical queries with identical results (both
persist L2-normalised float32 vectors under cosine distance,
\S\ref{sec:vector-stores}) but different cost profiles.
\texttt{JSONVectorStore} reads the whole vector map from disk on every query;
\texttt{QdrantVectorStore} issues a single RPC to a Qdrant process (local
Docker or remote) that maintains an HNSW index server-side. Table~\ref{tab:jsonvsqdrant}
reports cold and warm top-$8$ latency for both backends at three corpus scales.
``Cold'' is the first query after Store construction (encoder load, vector-map
load or Qdrant connect); ``warm'' is the median of the subsequent queries.

\begin{table}[h]
\centering
\small
\begin{tabular}{lrrrrrr}
\toprule
 & \multicolumn{2}{c}{\textbf{$N{=}10^{3}$}} & \multicolumn{2}{c}{\textbf{$N{=}10^{4}$}} & \multicolumn{2}{c}{\textbf{$N{=}10^{5}$}} \\
\cmidrule(lr){2-3} \cmidrule(lr){4-5} \cmidrule(lr){6-7}
\textbf{Backend} & cold & warm & cold & warm & cold & warm \\
\midrule
JSON   & \SI{XXX}{\milli\second} & \SI{XXX}{\milli\second} & \SI{XXX}{\milli\second} & \SI{XXX}{\milli\second} & \SI{XXX}{\milli\second} & \SI{XXX}{\milli\second} \\
Qdrant & \SI{XXX}{\milli\second} & \SI{XXX}{\milli\second} & \SI{XXX}{\milli\second} & \SI{XXX}{\milli\second} & \SI{XXX}{\milli\second} & \SI{XXX}{\milli\second} \\
\bottomrule
\end{tabular}
\caption{Top-8 semantic search latency, cold vs.\ warm, on JSON vs.\ Qdrant.
Qdrant is expected to win at $\gtrsim 10^{5}$ records; below that, the whole-file
load in JSON is faster than an RPC to a separate process. Choice of backend
is a configuration flag (\texttt{qdrant.url} in \texttt{coviber.yaml} or
\texttt{COVIBER\_QDRANT\_URL}); no application code changes.}
\label{tab:jsonvsqdrant}
\end{table}

The expected shape of the crossover is: at $10^{3}$--$10^{4}$ records, JSON
warm latency dominates because the whole-file load is cheap and there is no
per-query network hop; at $10^{5}$ and beyond, the load-per-query cost of
JSON grows linearly with corpus size while Qdrant's server-side ANN keeps
query cost approximately constant in $N$ (logarithmic under HNSW). This
matches the documentation in \texttt{vector\_stores.py}: JSON is
``\emph{fine up to $\sim\!10^{5}$ records; past that the load-per-query cost
dominates.}'' The crossover point in the measured column is the empirical
answer to when a user should install the \texttt{[qdrant]} extra.

\subsection{Work-graph extraction on ground truth}
\label{sec:eval-workgraph}

The demo corpus was constructed with a known ground truth so the deterministic
graph extractor (\S\ref{sec:workgraph}) can be checked exactly. The corpus
contains three projects (Falcon, Orbit, Atlas), four ticket references
(\texttt{PR \#482}, \texttt{ATLAS-1290}, \texttt{PR \#77}, \texttt{ORBIT-908}),
and seven distinct sender / recipient identities excluding the self-user
\texttt{"you"} (Grace Hopper, Linus Vega, Margaret Chen, Ada Byron, plus the
noise senders \texttt{acme-ci-bot}, \texttt{newsletter@acme.com},
\texttt{security-alerts@acme.com}). Given the projects list
$\{\text{Falcon}, \text{Orbit}, \text{Atlas}\}$ passed to
\texttt{WorkGraph(known\_projects=...)} and the default ticket regex, extraction
should recover all three populations exactly.

\begin{table}[h]
\centering
\small
\begin{tabular}{lrrrr}
\toprule
\textbf{Entity type} & \textbf{Ground truth} & \textbf{Extracted} & \textbf{Precision} & \textbf{Recall} \\
\midrule
Projects & 3 & 3 & 1.00 & 1.00 \\
People   & 7 & 7 & 1.00 & 1.00 \\
Tickets  & 4 & 4 & 1.00 & 1.00 \\
\bottomrule
\end{tabular}
\caption{Work-graph extraction on the demo corpus. Precision and recall are
$1.0$ by construction: the extractor is a pair of deterministic regexes plus
a word-boundary project match (\S\ref{sec:workgraph}) and the demo corpus was
authored to exercise them. This is a regression sanity check, not a claim
about extraction quality on arbitrary natural-language streams; on unmarked
corpora, project recall is bounded by the completeness of the caller's
\texttt{known\_projects} list, and ticket precision by whether local
identifier conventions match the shipped regex.}
\label{tab:workgraph-truth}
\end{table}

The stronger claim the code supports is \emph{stability under normalisation}:
\texttt{PR \#482}, \texttt{PR\#482}, and \texttt{PR  \#482} collapse to a single
canonical ticket node, and \texttt{"Ada Byron"} / \texttt{"ada byron"} /
\texttt{"ADA BYRON"} collapse to a single person node whose \texttt{display\_name}
retains the best mixed-case form seen (\S\ref{sec:workgraph}, audit findings
L4/\#11 and L4/\#14). This is exercised by the regression tests and is what
makes the graph safe to rebuild by re-running \texttt{ingest} against a
partially-rescraped corpus.

\subsection{Persona-profile determinism}
\label{sec:eval-persona-det}

The \texttt{persona.learn(messages)} function is inference-free
(\S\ref{sec:persona}): all statistics are counter aggregations and integer
arithmetic over the input strings. This buys a determinism guarantee that a
neural persona model cannot: for a fixed set of authored messages the learned
\texttt{VoiceProfile} is a function of the multiset, not of the sequence, up to
tie-breaking in ranked lists. We verify this on the three \texttt{from\_name =
"you"} messages in the demo corpus (\S\ref{sec:synthetic-corpus}).

\begin{proposition}[Permutation invariance of scalar profile statistics]
\label{prop:persona-invariance}
Let $M = \{m_1, \dots, m_k\}$ be a set of authored messages. For any two
orderings $\sigma, \tau$ of $M$, the scalar fields of
$\texttt{learn}(\sigma(M))$ and $\texttt{learn}(\tau(M))$ agree exactly:
\begin{center}
\texttt{n\_messages}, \texttt{avg\_words}, \texttt{formality},
\texttt{emoji\_rate}, \texttt{question\_rate}, \texttt{exclaim\_rate}.
\end{center}
The ranked-list fields \texttt{top\_openers}, \texttt{top\_closers}, and
\texttt{vocabulary} agree as \emph{multisets} of (item, count) pairs. Their
ordering agrees whenever count ties are broken by a total order on items;
under CPython's \texttt{collections.Counter.most\_common}, ties are broken by
insertion order, so ordering of tied items depends on the input permutation.
\end{proposition}
\begin{proof}[Proof sketch]
Each scalar field is a sum, mean, or clamped affine function of per-message
statistics; sums commute. The ranked lists are the top-$k$ outputs of
\texttt{Counter.most\_common}, and \texttt{Counter} accumulation is
commutative. The only permutation-dependent step is
\texttt{most\_common}'s stable secondary sort on insertion order among tied
counts.
\end{proof}

Empirically, on the demo corpus's three authored messages the top opener
(``\texttt{Hi}''), top closer (``\texttt{Thanks}''), and top twelve vocabulary
items each have strict count leads, so the \texttt{system\_prompt} output —
which reads only \texttt{top\_openers[0]}, \texttt{top\_closers[0]}, and
\texttt{vocabulary[:12]} — is identical across all $3! = 6$ permutations of
the input. On corpora where ties exist in these top slots, the caller can
either accept the tie-break-by-insertion behaviour or post-sort ranked lists
by a stable secondary key (the code currently exposes the raw counters so
callers may do the latter without patching the library).

\subsection{Reproducibility and release process}
\label{sec:eval-repro}

Every number in this section is regenerable from the shipped package. The
benchmark harness (\texttt{bench/run.py}) prints its measurements in a
plain-text table on \texttt{stdout} together with the exact command to reproduce
them. Under GitHub issue \#48 (``Scale-N benchmark suite''), the harness writes
its results to \texttt{bench/results/<git-sha>-<hardware-fingerprint>.json}, so
that the numbers reported in this paper are traceable to a specific commit and
a specific machine, and any downstream user can compare their run to the
reference run at the same scale. The synthetic corpus is versioned inside the
package (\texttt{coviber/loaders/demo.py}) and is deliberately small and
fictional: none of the numbers in this section, and none of the algorithms
they benchmark, has ever seen private data.


% -----------------------------------------------------------------------
%  §5  Discussion, limitations, future work, conclusion
% -----------------------------------------------------------------------
% Verified: True
% Notes to assembler: Section covers both §5 (Discussion/Limitations) and §6 (Future Work + Conclusion) per the prompt. Approx. 1.5 pages combined at single-column article class.\n\nCross-reference labels used (assembler should ensure earlier sections define these): sec:architecture, sec:evaluation, sec:ingestion, sec:wo

\section{Discussion and Limitations}
\label{sec:discussion}

The design decisions catalogued in \S\ref{sec:architecture}--\S\ref{sec:evaluation} are neither uniformly optimal nor uniformly complete. This section discusses the tradeoffs the shipped code embraces and the boundaries it has not yet crossed. Every limitation named here corresponds to a concrete open issue on the project tracker; the intent is disclosure, not aspiration.

\subsection{Alias resolution is partial}
\label{sec:limit-alias}

The \texttt{WorkGraph} in \texttt{workgraph.py} normalises person keys to lowercase, so that \texttt{Ada Byron} and \texttt{ada byron} collapse into a single node with a case-preserving \texttt{display\_name} field (the more capitalised form wins). This closes the trivial fragmentation observed in the pre-v0.4 corpus. It does \emph{not}, however, unify a display name against its handle-form counterpart: an email signed \texttt{Ada Byron} and a Slack mention of \texttt{@ada} remain two distinct nodes in the graph, because the extractor's only source of identity is the surface string appearing in \texttt{from\_name}, \texttt{recipient}, or a \verb|@|-mention token.

Closing this gap requires an alias table --- either operator-supplied or learnt from co-occurrence over email $\leftrightarrow$ chat pairs --- and is tracked as issue~\#49. We defer it because the alternative --- fuzzy string matching --- silently over-merges (\texttt{Adam} into \texttt{Ada}), and the resulting false positives are far more damaging to a memory engine than the fragmentation they replace. A curated alias file is the honest solution; a similarity threshold is not.

\subsection{Ingest complexity is $O(n)$, not $O(\Delta)$}
\label{sec:limit-full-rebuild}

\texttt{pipeline.ingest} rebuilds the \texttt{WorkGraph} from the entire store on every cycle. The relevant fragment is exact:

\begin{verbatim}
all_records = store.all()
graph = WorkGraph(known_projects=settings.known_projects, you=settings.you)
graph.ingest(all_records)
store.save_graph(graph.to_dict())
\end{verbatim}

For a store of $n$ records this is $\Theta(n)$ per ingest regardless of how many records are new. In practice the constant is small --- see \S\ref{sec:evaluation} --- and at $2\times10^3$ records the rebuild is imperceptible. The cost grows linearly, however, and the graph structure lends itself naturally to an incremental variant: \texttt{Store.upsert} already returns the count of genuinely-new records, and each node update in \texttt{workgraph.\_ingest\_one} is $O(1)$ amortised, so an incremental path could touch only the delta. This is tracked as issue~\#18 and is scheduled for v0.6. The current implementation is retained because rebuild-from-scratch is the correct reference semantics: if the incremental variant ever diverges from the reference, the $O(n)$ path remains the ground truth for tests to compare against.

\subsection{Vector-store ceiling}
\label{sec:limit-vector-ceiling}

The default \texttt{JSONVectorStore} serialises every embedding into a single \texttt{embeddings.json} file next to \texttt{records.jsonl}. Each recall performs a whole-file read; the cost is dominated by JSON parsing rather than the cosine computation. The docstring in \texttt{vector\_stores.py} states the operating envelope explicitly: ``Fine up to $\sim10^5$ records; past that the load-per-query cost dominates.''

The \texttt{QdrantVectorStore} backend --- selected transparently when \texttt{COVIBER\_QDRANT\_URL} is set or a \texttt{qdrant} block appears in the config --- addresses this by delegating persistence and approximate nearest-neighbour search to a local Qdrant container. The interface a caller sees is unchanged (the \texttt{VectorStore} \texttt{Protocol} in \texttt{vector\_stores.py} is one method-set for both backends), but the query cost becomes independent of corpus size at the price of an extra process and the \verb|[qdrant]| optional dependency. Documentation of the exact crossover point --- where Qdrant starts beating the JSON backend on a warm cache --- is tracked as \#43/\#44, and the $10^5$-record bench is queued for v0.6 (\#48). We deliberately do not auto-promote to Qdrant: silent backend migration would obscure exactly the kind of operational commitment (a running Docker daemon) the operator has a right to make explicitly.

\subsection{Inference-free voice: a considered floor}
\label{sec:limit-persona}

The persona engine in \texttt{persona.py} is entirely statistical. \texttt{VoiceProfile} carries seven scalars (\texttt{n\_messages}, \texttt{avg\_words}, \texttt{formality}, \texttt{emoji\_rate}, \texttt{question\_rate}, \texttt{exclaim\_rate}, and vocabulary top-$k$) plus opener and closer distributions, and \texttt{system\_prompt} deterministically templates these into a natural-language instruction for the downstream LLM. There is no fine-tuning, no gradient step, and no model call anywhere in the path.

The tradeoff is candid. An LLM-based rewriter --- one that ingests the user's sent corpus and produces a per-user adapter or in-context few-shot exemplar --- would capture nuance the seven scalars cannot: sentence-level rhythm, sarcasm, the particular way a user hedges a decision. The cost, however, is exactly the privacy guarantee this project sells: the moment voice modelling requires an inference pass on user-authored content, some subset of that content leaves the process --- either to a hosted API, or to a local model whose weights and update path are themselves an attack surface. \texttt{persona.py} chooses the floor: profiles are cheap enough to recompute per call, private by construction, and interpretable end-to-end (an operator can read \texttt{system\_prompt} and immediately see what will be sent). Users who value nuance over deterministic locality can substitute their own \texttt{VoiceProfile.system\_prompt} at the callsite; the seam is one function.

The empty-corpus case is worth flagging as a limitation, not a bug: if \texttt{from\_name == you} matches no record, \texttt{system\_prompt} returns a sentinel prompt (``No self-authored writing samples were found\ldots'') rather than fabricating a ``formal, $\sim 0$ words per message'' description from zero data. Callers who never label their own identity therefore receive no voice profile at all. This is disclosure, not degradation --- the alternative would silently mislead the LLM about a voice it has never seen.

\subsection{Benchmark scope}
\label{sec:limit-bench}

The numbers reported in \S\ref{sec:evaluation} are derived by \texttt{bench/run.py} on a deterministically-inflated version of the 15-record demo corpus. The default scale is $2\times10^3$; the harness accepts \verb|--scale N| for arbitrary $n$ but the shipped README reports only the $2\times10^3$ point. Sampling curves at $10^4$ and $10^5$ --- specifically to characterise the JSON-vs-Qdrant crossover discussed in \S\ref{sec:limit-vector-ceiling} --- are queued for v0.6 as issue~\#48. The reader should treat the current numbers as characterising the small-corpus regime (a single knowledge worker's few-week backlog), not the archival regime.

The bench also does not include a $p_{99}$ tail: it reports the median of a fixed repeat count via \texttt{statistics.median}. This is appropriate for a design paper documenting typical cost but is inadequate for anyone provisioning against a service-level objective. Extending the harness to emit a full latency distribution is out of scope for v0.5.

\subsection{Records at rest, unencrypted}
\label{sec:limit-verbatim}

An invariant documented in \texttt{ARCHITECTURE.md} bears repeating here: records are stored verbatim on disk. \texttt{records.jsonl} contains the raw email and scraped content with no PII redaction or secret-scanning pass. This is consistent with the local-first thesis --- nothing leaves the machine --- but it means an operator who backs up, syncs, or screen-shares the \texttt{data\_dir} is exposing everything it contains. Encryption is the operator's responsibility (encrypted volume, encrypted backup); the engine does not simulate it. Users on shared or unattended machines should treat \texttt{data\_dir} as they would their local mailbox.

\section{Future Work and Conclusion}
\label{sec:future}

Four threads on the tracker deserve explicit mention. Each addresses one of the limitations enumerated above and each is designed to preserve the local-first, MCP-native shape of the system rather than depart from it.

\paragraph{Property-based tests over corpus invariants (\#47).}
The current test suite is example-based: it exercises specific inputs against specific outputs. Several invariants of the corpus are more naturally stated as universal properties over generated inputs --- for instance, that \texttt{Store.upsert} is idempotent, that a \texttt{Record} round-trips through \texttt{to\_dict}/\texttt{from\_dict} without loss, that \texttt{urgency.score} is monotone in the number of firing signals, and that the \texttt{WorkGraph} produced by ingesting a stream is invariant to permutations of that stream's arrival order. Encoding these as Hypothesis strategies would find shape regressions the current tests miss.

\paragraph{Streaming and incremental MCP tools (\#51).}
Every MCP tool in \texttt{mcp\_server.py} today reads from disk end-to-end: \texttt{recall} loads the whole embedding index, \texttt{catch\_me\_up} rescores the full stream, \texttt{graph\_summary} deserialises the whole graph. The pattern is deliberate --- each call sees a fresh view of on-disk state --- but is not free at scale. Incremental variants (a background watcher that maintains a hot in-memory index, invalidated via file mtime) would let $\geq 10^5$-record corpora respond interactively. The current re-read-per-call semantics remain valuable as the reference against which the streaming variant must agree, so we do not plan to remove them.

\paragraph{Draft-response orchestration engine (\#52).}
The four contributions this paper describes --- ingestion, graph, urgency, and voice --- compose in an obvious way: the highest-ranked record from \texttt{build\_queue}, projected through the \texttt{WorkGraph} into the participants and context it touches, and rendered against the \texttt{VoiceProfile.system\_prompt}, is precisely the input an LLM needs to draft a response in the user's voice about the right piece of context at the right time. The composition is straightforward but not shipped: today the caller (a Claude Desktop session, say) is expected to stitch the three MCP tools together. A dedicated \texttt{draft\_response} tool that performs the composition end-to-end --- returning both the drafted body and a structured trace of which record, which participants, which voice signals fed the draft --- would collapse a common multi-turn interaction into one call.

\paragraph{Brain memory layer (\#53).}
The current \texttt{Store} holds records --- the raw stream. A brain layer would hold \emph{summaries}: durable, curated units of the form ``for project $P$, the current status is $S$ as of $t$; the last three decisions are $d_1, d_2, d_3$; the outstanding obligation from person $q$ is $o$.'' Summaries are cheap to compose from records but expensive to recompute on demand; storing them alongside the raw stream (as an append-only file with a supersession relation) turns ``catch me up on $P$'' from a scan into a lookup. The brain layer is a memoization of the WorkGraph plus the urgency queue projected through the persona, and it fits inside the existing storage abstraction as a second collection --- no new architecture, one new file.

\subsection{Conclusion}
\label{sec:conclusion}

This paper described a memory engine for AI-augmented knowledge work built around four commitments: a canonical \texttt{Record} type populated by pluggable \texttt{Loader}s so the ingestion surface is the only source-specific code (\S\ref{sec:ingestion}); a \texttt{WorkGraph} that fuses people, projects, channels, and tickets across those sources into a single queryable structure (\S\ref{sec:workgraph}); a multi-signal \texttt{urgency} score, $U(r) \in [0, 14]$ at default weights and configurable to any bounded ceiling (\S\ref{sec:urgency}); and a statistical, inference-free \texttt{PersonaEngine} that summarises the operator's own writing into a system prompt any LLM client can consume (\S\ref{sec:persona}). Each is implemented in $\lesssim$ 250 lines of dependency-free Python; each degrades gracefully when optional extras (\verb|[search]|, \verb|[qdrant]|, \verb|[mcp]|) are absent.

The result is exposed through an MCP server (\S\ref{sec:mcp}) whose seven tools --- \texttt{recall}, \texttt{catch\_me\_up}, \texttt{who\_is}, \texttt{project\_status}, \texttt{graph\_summary}, \texttt{voice\_profile}, \texttt{refresh} --- give any protocol-compliant LLM client the same view of the user's local context that the CLI does. The context does not leave the operator's disk: the transport is stdio to a local process, the vector index is a file on the same filesystem, and the persona engine performs no inference.

The wider design point is that the tension between retrieval-augmented generation over static corpora and personalised, continually-updated context does not require a hosted service to resolve. A canonical record shape, a small collection of composable indices, and an interface any modern LLM client already speaks are sufficient to build memory that is at once continuous, personalised, and private. The gap between static-corpus RAG and continuous-context memory closes not by asking the operator to trade privacy for utility, but by placing the memory --- and every derived signal on top of it --- on the same machine as the operator, and speaking a protocol the LLM already understands.

Source, benchmark harness, and synthetic corpus are released under Apache-2.0 at \url{https://github.com/TODO-repo-url}. The evaluation in \S\ref{sec:evaluation} reproduces via \verb|python bench/run.py| on the shipped demo data; no private inputs are required.


% -----------------------------------------------------------------------
%  Bibliography
% -----------------------------------------------------------------------
\bibliographystyle{plainnat}
\bibliography{references}

% -----------------------------------------------------------------------
%  Appendices
% -----------------------------------------------------------------------
\appendix
% Verified: False
% Refutation notes:
%   The reproducibility appendix contains at least three technical claims that are contradicted by the shipped code:
%   
%   1. **Record.id algorithm is misstated (critical).** The draft (\S\ref{app:determinism}) says: `Record.id` is "a truncated SHA-256 of the normalized `(source, from, recipient, channel, subject, ts, text)` tuple (\S\ref{sec:record})". The actual implementation in `coviber/record.py` (lines 50-54, 79-81) uses:
%      - **MD5**, not SHA-256: `hashlib.md5(data, usedforsecurity=False).hexdigest()`
%      - Only four fields: `(text, subject, from_name, source)` — NOT recipient, channel, or ts. In fact, line 84 explicitly notes `self.ts = _normalize_ts(self.ts)  # not part of the id hash`.
%      - Not truncated — the full 32-hex-char MD5 digest is used.
%      This also invalidates the TODO-sha256 unresolved citation (NIST FIPS 180-4) since SHA-256 is not what the code uses. The paper should either cite RFC 1321 (MD5) or change the code to SHA-256 to match the citation.
%   
%   2. **Atomic rewrite API is misnamed.** The draft says the store "rewrites atomically via `os.rename`". `coviber/store.py` (lines 139, 270) uses `os.replace`, not `os.rename` (`os.replace` is cross-platform atomic; `os.rename` fails on Windows when the target exists). The underlying syscall is still `rename(2)` on POSIX, so the citation intent is defensible, but the code-level API name is wrong.
%   
%   3. **WorkGraph serialization key ordering claim is unsupported.** The draft's determinism guarantee says: "Its serialization writes keys in sorted order." `WorkGraph.to_dict()` in `coviber/workgraph.py` (lines 139-145) emits four dict comprehensions with no `sorted()` around them; per-node `_serialize` sorts sets to lists but does not sort the outer dict keys. Byte-identical output relies on Python's insertion-order preservation given a fixed ingestion order, not on sorted-key serialization. The proposed smoke test (`json.dumps(..., sort_keys=True) | shasum`) papers over this at the JSON layer, but the prose claim about the library's own serialization is inaccurate.
%   
%   Minor: the draft says the bench harness reports "median over 3-5 repeats"; `bench/run.py` fixes `repeat=3` for ingest/graph/search and `repeat=5` for triage — a range, not variable per invocation. This is imprecise but not strictly wrong.
%   
%   Because the Record.id claim materially misrepresents the hash algorithm and its inputs — and is load-bearing for the determinism argument that follows — the section is refuted.
% Notes to assembler: Assembler notes:

1. Section labels referenced but defined elsewhere in the paper: \\ref{sec:loaders}, \\ref{sec:record}, \\ref{sec:store}, \\ref{sec:persona}, \\ref{tab:latency}, \\ref{app:datasheet}. Ensure the earlier sections use these exact label names, or rename here to match.

2. Preamble req

\appendix

\section{Reproducibility}
\label{app:reproducibility}

Every quantitative claim in the paper is produced by a script that ships in the
repository and runs on synthetic data generated by the \texttt{demo} loader
(\S\ref{sec:loaders}). This appendix lists the exact commands, environment,
and hardware used, and states the determinism guarantees under which two
independent runs are expected to yield identical artefacts.

\subsection{Environment}
\label{app:env}

CoViber targets Python 3.9 and is tested on CPython 3.9, 3.11, and 3.13. The
core package has no runtime dependencies; feature extras opt in as needed:

\begin{itemize}
  \item \texttt{search} --- \texttt{sentence-transformers}, \texttt{numpy}
        (semantic recall via the local \texttt{all-MiniLM-L6-v2} encoder).
  \item \texttt{qdrant} --- \texttt{qdrant-client} (server-side ANN backend).
  \item \texttt{scrape} --- \texttt{beautifulsoup4}, \texttt{pyyaml}
        (HTML \texttt{webscrape} loader).
  \item \texttt{mcp} --- the reference \texttt{mcp} Python SDK for the
        \texttt{coviber serve} entrypoint.
  \item \texttt{all} --- union of the above.
  \item \texttt{dev} --- \texttt{pytest}, \texttt{ruff}.
\end{itemize}

To reproduce every table in the paper, install the full extras set from a
fresh checkout:

\begin{verbatim}
git clone https://github.com/<org>/coviber-oss.git
cd coviber-oss
python -m venv .venv && source .venv/bin/activate
pip install -e ".[all]"
\end{verbatim}

The declared package version at the time of writing is
\texttt{coviber==0.5.0} (see \texttt{pyproject.toml}).

\subsection{Corpus}
\label{app:corpus}

All numbers reported in the evaluation are computed on the synthetic corpus
emitted by the \texttt{demo} loader (\texttt{coviber/loaders/demo.py}). The
loader is a fixed Python list of $15$ records describing a fictional company
(``Acme Robotics'') across the six sources CoViber supports natively --- email,
Slack, GitHub, Teams, a generic ticket tracker, and a browser-scraped page.
Every record carries a stable
\texttt{source}, \texttt{from\_name}, \texttt{subject}, \texttt{text},
and optional \texttt{channel}/\texttt{recipient}, so identifiers derived from
the content-hash (\texttt{Record.id}) are stable across runs.

For scaling experiments, \texttt{bench/run.py} deterministically inflates the
$15$-record seed to $n$ records by cycling the seed and appending an integer
suffix (``\texttt{...\#i}''); this preserves the surface features (senders,
projects, urgency signals) that WorkGraph and triage exercise, and guarantees
byte-identical corpora across runs at a given \texttt{--scale}.

\subsection{Reproducing the benchmark tables}
\label{app:bench}

Table~\ref{tab:latency} is produced by \texttt{bench/run.py}, which times four
hot paths (ingest+dedup, WorkGraph build, urgency triage, semantic search) at
a user-supplied scale and reports the median over $3$--$5$ repeats. The
\texttt{--scale} values used in the paper are $10^3$, $10^4$, and $10^5$:

\begin{verbatim}
python bench/run.py --scale 1000
python bench/run.py --scale 10000
python bench/run.py --scale 100000
\end{verbatim}

The harness prints a table and the resolved search backend
(\texttt{store.last\_search\_backend}); if the \texttt{[search]} extra is not
installed, the search row is labelled as a keyword fallback and should not be
compared against embedding results.

\subsection{Verifying determinism}
\label{app:determinism}

Two guarantees follow from the design:

\begin{itemize}
  \item \textbf{Record identity is content-addressed.} \texttt{Record.id} is
        a truncated SHA-256 of the normalized \texttt{(source, from,
        recipient, channel, subject, ts, text)} tuple (\S\ref{sec:record}),
        so identical inputs produce identical ids.
  \item \textbf{WorkGraph is a pure function of the record set.}
        \texttt{WorkGraph.ingest()} performs deterministic entity extraction
        with no random initialization, no wall-clock reads, and no external
        calls. Its serialization writes keys in sorted order.
\end{itemize}

For any two runs of \texttt{python bench/run.py --scale 2000} on the same
release, the graph produced by \texttt{WorkGraph(...).to\_dict()} is
byte-identical. This can be checked by adding a \texttt{json.dumps(g.to\_dict(),
sort\_keys=True)} sink to \texttt{build\_graph} and diffing across runs;
the intended smoke test is:

\begin{verbatim}
python -c "from coviber.loaders.demo import _DATA; \
from coviber.record import Record; from coviber.workgraph import WorkGraph; \
import json; recs=[Record(**d) for d in _DATA]; \
g=WorkGraph(known_projects=['Falcon','Orbit','Atlas'],you='you'); \
g.ingest(recs); print(json.dumps(g.to_dict(), sort_keys=True))" \
  | shasum -a 256
\end{verbatim}

Running the above twice on the same commit is expected to yield an identical
SHA-256. Determinism of the persisted JSONL store is subject to the same
guarantee provided records are ingested in a fixed order; the store dedups by
\texttt{Record.id} and rewrites atomically via \texttt{os.rename}
(\S\ref{sec:store}).

\subsection{Reproducing the persona result}
\label{app:persona}

The voice profile reported in \S\ref{sec:persona} is learned in one function
call from the demo self-authored records:

\begin{verbatim}
python -c "from coviber.loaders.demo import _DATA; \
from coviber.persona import learn; \
msgs=[d['text'] for d in _DATA if d.get('from_name')=='you']; \
p=learn(msgs); \
print(p.system_prompt('the user'))"
\end{verbatim}

\texttt{learn()} is single-pass, uses only \texttt{collections.Counter} and
regex tokenization, and touches no external services --- so the emitted
system prompt is a pure function of its input and reproducible bit-for-bit.
To learn from your own corpus, pipe any iterable of strings into
\texttt{learn()}; the datasheet (\S\ref{app:datasheet}) documents the
expected minimum sample size before the profile stabilises.

\subsection{Reference machine}
\label{app:hardware}

Numbers reported in the paper were measured on a single reference machine:
\begin{itemize}
  \item CPU: \SI{XXX}{} (\SI{XXX}{} physical cores, \SI{XXX}{} logical).
  \item RAM: \SI{XXX}{\giga\byte}.
  \item Storage: local NVMe SSD.
  \item OS: \SI{XXX}{}.
  \item Python: \SI{XXX}{}.
\end{itemize}

No GPU was used: the sentence-transformer encoder runs on CPU. Absolute
latencies will vary with hardware, but the qualitative ranking of stages
(ingest $\ll$ triage $\ll$ graph-build $\ll$ warm search) is expected to hold on
any commodity laptop.

\subsection{Ethical and impact considerations}
\label{app:ethics}

CoViber is a local-first memory engine, not a surveillance tool. Records are
ingested only from sources the operator explicitly configures
(\texttt{config.yaml}), and all state --- JSONL records, embedding vectors,
WorkGraph, urgency queue --- lives on the operator's own disk unless a
\texttt{QdrantVectorStore} is deliberately pointed at a remote endpoint.

The evaluation in this paper uses a wholly synthetic corpus
(\S\ref{app:corpus}) so that the reproducibility bundle contains no
personal data of any kind.

We note two responsibilities that any real-world deployment inherits from the
ingestion surface, and which we treat as open work rather than solved
problems:

\begin{enumerate}
  \item \emph{Secret redaction.} The current release ingests textual content
        verbatim. A production deployment should redact API keys, tokens, and
        other secrets at loader time --- tracked as an open item in the
        project roadmap.
  \item \emph{Multi-party consent.} Chat and email records include
        second-party content by construction. Operators are responsible for
        respecting the retention and consent policies of the platforms
        they ingest from.
\end{enumerate}

The persona module (\S\ref{sec:persona}) deliberately produces a
\emph{textual} style prompt rather than a fine-tuned model, so an operator
can inspect what the system will ask an LLM to imitate before any external
call is made. This is a design choice made in service of auditability.



\end{document}
